% future/future.tex
% mainfile: ../perfbook.tex
% SPDX-License-Identifier: CC-BY-SA-3.0

\QuickQuizChapter{chp:Conflicting Visions of the Future}{对未来的不同愿景}{qqzfuture}
%
\Epigraph{预测是非常困难的，尤其是关于未来的预测。}
	 {Niels Bohr}

本章介绍了关于并行编程未来的一些相互矛盾的愿景。
目前尚不清楚这些愿景中哪些会成为现实，事实上，也不清楚其中是否有任何一个会实现。
然而它们仍然很重要，因为每种愿景都有其忠实的拥护者，如果有足够多的人对某件事
抱有足够狂热的信念，你就需要面对该事物以其对拥护者的思想、言论和行动的影响形式
而存在的事实。
此外，这些愿景中的一个或多个确实会成为现实。
但大多数是不靠谱的。
能分辨出哪个是哪个，你就发财了~\cite{KeithRSpitz1977}！

因此，以下各节概述了transactional memory、hardware transactional memory、
回归测试中的形式化验证以及并行functional programming。
但首先，让我们看一个来自21世纪初的关于预测的警示故事。

% future/cpu.tex
% mainfile: ../perfbook.tex
% SPDX-License-Identifier: CC-BY-SA-3.0

\section{CPU技术的未来已今非昔比}
\label{sec:future:The Future of CPU Technology Ain't What it Used to Be}
%
\epigraph{辉煌的未来已成过去。}{David Maraniss}

根据多年的经验，回首过去的岁月总是那么简单而纯真。
21世纪初，最大的无知表现在：大多数人对\IXr{Moore's Law}
即将无法继续带来传统CPU时钟频率提升这一事实浑然不知。
当然，偶尔也有关于技术极限的警告，但这样的警告已经发出了数十年。
怀着这样的心态，请考虑以下几种情景：

\begin{figure}
\centering
\resizebox{3in}{!}{\includegraphics{cartoons/r-2014-CPU-future-uniprocessor-uber-alles}}
\caption{单处理器至上}
\ContributedBy{fig:future:Uniprocessor \"Uber Alles}{Melissa Broussard}
\end{figure}

\begin{figure}
\centering
\resizebox{2.6in}{!}{\includegraphics{cartoons/r-2014-CPU-Future-Multithreaded-Mania}}
\caption{多线程狂热}
\ContributedBy{fig:future:Multithreaded Mania}{Melissa Broussard}
\end{figure}

\begin{figure}
\centering
\resizebox{2.5in}{!}{\includegraphics{cartoons/r-2014-CPU-Future-More-of-the-Same}}
\caption{一切照旧}
\ContributedBy{fig:future:More of the Same}{Melissa Broussard}
\end{figure}

\begin{figure}
\centering
\resizebox{3in}{!}{\includegraphics{cartoons/r-2014-CPU-Future-Crash-dummies}}
\caption{撞向内存墙的碰撞试验假人}
\ContributedBy{fig:future:Crash Dummies Slamming into the Memory Wall}{Melissa Broussard}
\end{figure}

\begin{figure}
\centering
\resizebox{3in}{!}{\includegraphics{cartoons/r-2021-CPU-future-astounding-accelerator}}
\caption{令人惊叹的加速器}
\ContributedBy{fig:future:Astounding Accelerators}{Melissa Broussard, remixed}
\end{figure}

\begin{enumerate}
\item	单处理器至上
	(\cref{fig:future:Uniprocessor \"Uber Alles})，
\item	多线程狂热
	(\cref{fig:future:Multithreaded Mania})，
\item	一切照旧
	(\cref{fig:future:More of the Same})，以及
\item	撞向内存墙的碰撞试验假人
	(\cref{fig:future:Crash Dummies Slamming into the Memory Wall})。
\item	令人惊叹的加速器
	(\cref{fig:future:Astounding Accelerators})。
\end{enumerate}

以下各节分别介绍每种情景。

\subsection{单处理器至上}
\label{sec:future:Uniprocessor \"Uber Alles}

正如2004年所说~\cite{PaulEdwardMcKenneyPhD}：

\begin{quote}
	在这一情景中，\IXaltr{Moore's-Law}{Moore's Law}带来的CPU时钟频率提升
	与水平扩展计算的持续进步相结合，使得SMP系统变得无关紧要。
	因此，这一情景被称为``单处理器至上''，
	字面意思是单处理器高于一切。

	这些单处理器系统将只受到指令开销的影响，因为\IXpl{memory barrier}、
	缓存抖动和竞争不会影响单CPU系统。
	在这一情景中，RCU仅对小众应用有用，例如
	与\IXacrpl{nmi}交互。
	目前还不清楚一个缺少RCU的操作系统是否会觉得有必要采用它，
	尽管已经实现了RCU的操作系统可能会继续使用。

	然而，多线程CPU的最新进展似乎表明
	这一情景是极不可能的。
\end{quote}

确实极不可能！
但更广泛的软件社区不愿接受他们需要拥抱并行性这一事实，
还需要一段时日，这个社区才会得出结论：
\IXaltr{Moore's-Law}{Moore's Law}带来的CPU时钟频率提升这顿``免费午餐''
确实已经彻底结束了。
永远不要忘记：
信念是一种情感，而不是理性技术思考的结果！

\subsection{多线程狂热}
\label{sec:future:Multithreaded Mania}

同样来自2004年~\cite{PaulEdwardMcKenneyPhD}：

\begin{quote}
	单处理器至上有一种不那么极端的变体，即具有硬件多线程的单处理器，
	事实上，多线程CPU现在已经是许多台式和笔记本电脑系统的标配。
	最激进的多线程CPU共享所有级别的缓存层次结构，
	从而消除了CPU间的\IXh{memory}{latency}，
	相应地也大大降低了传统同步机制的性能损失。
	然而，多线程CPU仍然会因竞争和内存屏障（memory barrier）导致的
	流水线停顿而产生开销。
	此外，由于所有硬件线程共享所有级别的缓存，
	单个硬件线程可用的缓存只是等效单线程CPU上缓存的一小部分，
	这可能会降低大量使用缓存的应用程序的性能。
	还存在一种可能性，即有限的可用缓存将导致基于RCU的算法
	因其 grace period 引起的额外内存消耗而遭受性能损失。
	调查这种可能性是未来的工作。

	然而，为了避免这种性能下降，许多多线程CPU和多CPU芯片
	至少在某些缓存级别上按每个硬件线程进行分区。
	这增加了每个硬件线程可用的缓存量，
	但重新引入了硬件线程间传递缓存行的内存延迟。
\end{quote}

最终，我们都知道这个故事是如何发展的：
单个芯片上有多个多线程核心，插在单个插槽中，
对较少的活跃线程数有不同程度的优化。
问题是，未来的共享内存系统是否总是能装入单个插槽中。

\subsection{一切照旧}
\label{sec:meas:More of the Same}

再次来自2004年~\cite{PaulEdwardMcKenneyPhD}：

\begin{quote}
	一切照旧情景假设，内存延迟比率将大致保持在今天的水平。

	这一情景实际上代表了一种变化，因为要保持不变，
	互连性能必须跟上\IXaltr{Moore's-Law}{Moore's Law}
	带来的核心CPU性能提升。
	在这一情景中，由于流水线停顿、内存延迟和竞争引起的开销
	仍然十分显著，RCU仍然保持着它今天所享有的高适用性。
\end{quote}

而变化正是\IXr{Moore's Law}仍在提供的不断提升的集成度。
但从长远来看，到底是什么？
在每个芯片上放置更多的CPU？
还是更多的I/O、缓存和内存？

服务器似乎在选择前者，而片上系统（SoC）的嵌入式系统则继续选择后者。

\subsection{撞向内存墙的碰撞试验假人}
\label{sec:future:Crash Dummies Slamming into the Memory Wall}

\begin{figure}
\centering
\epsfxsize=3in
\epsfbox{future/latencytrend}
% from Ph.D. thesis: related/latencytrend.eps
\caption{Sequent计算机每次本地内存引用的指令数}
\label{fig:future:Instructions per Local Memory Reference for Sequent Computers}
\end{figure}

\begin{figure}
\centering
\epsfxsize=3in
\epsfbox{future/be-lb-n4-rf-all}
% from Ph.D. thesis: an/plots/be-lb-n4-rf-all.eps
\caption{盈亏平衡点 vs.\@ $r$，$\lambda$ 较大，四个CPU}
\label{fig:future:Breakevens vs. r; lambda Large; Four CPUs}
\end{figure}

\begin{figure}
\centering
\epsfxsize=3in
\epsfbox{future/be-lw-n4-rf-all}
% from Ph.D. thesis: an/plots/be-lw-n4-rf-all.eps
\caption{盈亏平衡点 vs.\@ $r$，$\lambda$ 较小，四个CPU}
\label{fig:future:Breakevens vs. r; Worst-Case lambda; Four CPUs}
\end{figure}

2004年的又一段引言~\cite{PaulEdwardMcKenneyPhD}：

\begin{quote}
	如果
	\cref{fig:future:Instructions per Local Memory Reference for Sequent Computers}
	中所示的内存延迟趋势继续下去，那么内存延迟将继续相对于
	指令执行开销增长。
	像Linux这样大量使用RCU的系统将发现进一步使用RCU是有利可图的，
	如\cref{fig:future:Breakevens vs. r; lambda Large; Four CPUs}所示。
	从该图可以看出，如果RCU被大量使用，增加的内存延迟比率
	使RCU相对于其他同步机制获得越来越大的优势。
	相比之下，很少使用RCU的系统将需要越来越高的读密集度
	才能使RCU的使用有所回报，
	如\cref{fig:future:Breakevens vs. r; Worst-Case lambda; Four CPUs}所示。
	从该图可以看出，如果RCU使用较少，增加的内存延迟比率
	使RCU与其他同步机制相比处于越来越不利的地位。
	由于在高负载下Linux已被观察到每个\IX{grace period}
	超过1,600个回调~\cite{Sarma04c}，
	可以安全地说Linux属于前一类。
\end{quote}

一方面，这段话未能预见到缓存预热的问题，即RCU
在承受大量高强度更新写负载时可能遭受的性能损失，
大概是因为当时看来RCU不太可能用于这种工作负载。
实际上，\co{SLAB_TYPESAFE_BY_RCU}已在许多实例中投入使用，
在这些情况下否则会出现缓存预热问题，sequence locking也是如此。
另一方面，这段话也未能预见到RCU会被用于减少调度延迟或用于安全目的。

本书生成的大部分数据是在一个八插槽系统上收集的，
每个插槽有28个核心，每个核心有两个硬件线程，
总共448个硬件线程。
空闲系统的内存延迟小于一微秒，
与2004年同等规模系统的延迟不相上下。
有些人声称这些延迟接近一微秒仅仅是因为x86 CPU系列
相对较强的内存排序，但可能还需要一段时间
才能解决这一特定争论。

\subsection{令人惊叹的加速器}
\label{sec:future:Astounding Accelerators}

硬件加速器的潜力在2004年不像在2021年那样清晰，
因此本节没有引言。
然而，2020年11月的Top 500列表~\cite{Top500}包含大量加速器，
因此可以说本节是对现在而非未来的展望。
前面几节中的大多数也可以这样说。

硬件加速器正被用于许多其他用途，包括
加密、压缩和机器学习。

简而言之，还是要小心进行预测，包括本章其余部分的预测。

\IfTwoColumn{}{\FloatBarrier}
% future/tm.tex
% mainfile: ../perfbook.tex
% SPDX-License-Identifier: CC-BY-SA-3.0

\section{事务内存}
\label{sec:future:Transactional Memory}
\epigraph{一切应该尽可能简单，但不能过于简单。}
	 {Albert Einstein，经由Louis Zukofsky}


在数据库之外使用事务的想法可以追溯到几十年
前~\cite{DBLomet1977SIGSOFT,Knight:1986:AMF:319838.319854,Herlihy93a}，
数据库事务和非数据库事务之间的关键区别在于，非数据库事务
去掉了定义数据库事务的``ACID''\footnote{
	原子性（Atomicity）、一致性（Consistency）、隔离性（Isolation）
	和持久性（Durability）。}
属性中的``D''。
在硬件中支持基于内存的事务，即``transactional memory''
（TM）的想法出现得更
晚~\cite{Herlihy93a}，但不幸的是，尽管也有一些类似的提案被提
出~\cite{JMStone93}，商用硬件对此类事务的支持并未立即出现。
不久之后，Shavit和Touitou提出了一种纯软件的
transactional memory (STM) 实现，它能够在商用硬件上运行，
尽管存在内存排序问题~\cite{Shavit95}。
这一提案沉寂了多年，也许是因为研究社区的注意力被\IXacrl{nbs}
（见\cref{sec:advsync:Non-Blocking Synchronization}）所吸引。

但到了世纪之交，TM开始受到更多
关注~\cite{Martinez01a,Rajwar01a}，到了这个十年的中期，
人们的兴趣只能用``白热化''来形
容~\cite{MauriceHerlihy2005-TM-manifesto.pldi,
DanGrossman2007TMGCAnalogy}，只有少数谨慎的
声音~\cite{Blundell2005DebunkTM,McKenney2007PLOSTM}。

TM背后的基本思想是原子地执行一段代码，使得其他线程（thread）看不到中间状态。
因此，TM的语义可以简单地通过将每个事务替换为一个可递归获取的
全局锁的获取和释放来实现，尽管性能和可扩展性会极其糟糕。
TM实现（无论是硬件还是软件）中固有的大部分复杂性在于
高效地检测并发事务何时可以安全地并行运行。
由于这种检测是动态进行的，冲突的事务可能会被中止或``回滚''，
而且在某些实现中，这种失败模式对程序员是可见的。

由于事务回滚的可能性随着事务大小的减小而降低，
TM对于小型基于内存的操作可能变得相当有吸引力，
例如用于栈、队列、哈希表（hash table）和搜索树的链表操作。
然而，目前对于大型事务，特别是包含I/O和进程创建等
非内存操作的事务，要证明其适用性则困难得多。
以下各节审视了``Transactional Memory
Everywhere''~\cite{PaulEMcKenney2009TMeverywhere}
这一宏大愿景当前面临的挑战。
\Cref{sec:future:Outside World}考察了与外部世界交互所面临的挑战，
\cref{sec:future:Process Modification}探讨了与进程修改原语的交互，
\cref{sec:future:Synchronization}探索了与其他同步原语的交互，
最后\cref{sec:future:Discussion}以一些讨论作为结尾。

\subsection{外部世界}
\label{sec:future:Outside World}

用\ppl{Donald}{Knuth}的智慧之言来说：

\begin{quote}
	许多计算机用户觉得输入和输出实际上不是``真正的编程''的一部分，
	它们仅仅是为了将信息输入和输出机器而（不幸地）
	必须做的事情。
\end{quote}

无论我们是否相信输入和输出是``真正的编程''，
事实是软件绝对必须与外部世界打交道。
因此，本节评析transactional memory的外部世界能力，
重点关注I/O操作、时间延迟和持久存储。

而这些与外部世界的交互正是TM无条件可组合性声明
破灭的基石。
是的，你可以组合事务，但前提是没有介入的TM不友好操作。
正如你可以组合基于锁的临界区（critical section），但前提是这样做不会引入死锁（deadlock）。

最终，在实际代码中TM是否真的比锁提供更好的可组合性，
这一点并不清楚。

\subsubsection{I/O操作}
\label{sec:future:I/O Operations}

可以在基于锁的临界区内执行I/O操作，
在持有\IX{hazard pointer}时，在sequence-locking读侧临界区内，
以及在userspace-RCU读侧临界区内执行I/O操作，
如果需要的话，甚至可以同时在所有这些情况下执行。
当你尝试从事务内部执行I/O操作时会发生什么？

根本问题在于事务可能会被回滚，例如由于冲突。
粗略地说，这要求给定事务内的所有操作都是可撤销的，
即执行两次操作与执行一次具有相同的效果。
不幸的是，I/O通常是典型的不可撤销操作，
这使得将通用I/O操作包含在事务中变得困难。
事实上，通用I/O是不可撤销的：
一旦你按下了发射核弹头的那个众所周知的按钮，
就再也无法回头了。

以下是在事务中处理I/O的一些选项：

\begin{enumerate}
\item	将事务内的I/O限制为使用内存缓冲区的缓冲I/O。
	这些缓冲区可以像包含任何其他内存位置一样被包含在事务中。
	这似乎是首选机制，而且在许多常见的流I/O和
	大容量存储I/O场景中确实工作得很好。
	然而，在多个面向记录的输出流从多个进程合并到单个文件的情况下，
	需要特殊处理，例如使用\co{fopen()}的``a+''选项
	或\co{open()}的\co{O_APPEND}标志时。
	此外，如下一节所述，常见的网络操作无法通过缓冲来处理。
\item	禁止事务内的I/O，使得任何执行I/O操作的尝试都会中止
	包含它的事务（也许还有多个嵌套事务）。
	这种方法似乎是非缓冲I/O的传统TM处理方式，
	但要求TM与容忍I/O的其他同步原语互操作。
\item	禁止事务内的I/O，但借助编译器来强制执行这一禁令。
\item	在任何给定时间只允许一个特殊的
	\emph{不可撤销}事务~\cite{SpearMichaelScott2008InevitableSTM}
	进行，
	从而允许不可撤销事务包含I/O操作。\footnote{
		在早期文献中，不可撤销事务被称为
		\emph{inevitable} 事务。}
	这在一般情况下可行，但严重限制了I/O操作的可扩展性和性能。
	鉴于可扩展性和性能是并行性的首要目标，
	这种方法的通用性似乎有些自我限制。
	更糟糕的是，使用不可撤销性来容忍I/O操作似乎极大地限制了
	手动事务中止操作的使用。\footnote{
		这一困难由Michael Factor指出。
		要理解这个问题，请思考一下当尝试在事务执行了
		不可撤销操作之后中止该事务时，TM应该做什么。}
	最后，如果有一个不可撤销事务正在操作某个给定的数据项，
	那么操作同一数据项的任何其他事务都不能具有非阻塞语义。
\item	创建新的硬件和协议，使I/O操作能够被纳入事务基础设施中。
	在输入操作的情况下，硬件需要正确预测操作的结果，
	并在预测失败时中止事务。
\end{enumerate}

I/O操作是TM众所周知的弱点，而且在事务中支持I/O的问题
是否有合理的通用解决方案尚不清楚，至少如果``合理''
要包含可用的性能和可扩展性的话。
尽管如此，持续对这一问题的关注和投入很可能会带来更多进展。

\subsubsection{RPC操作}
\label{sec:future:RPC Operations}

可以在基于锁的临界区内执行RPC，在持有hazard pointer时，
在sequence-locking读侧临界区内，以及在userspace-RCU读侧
临界区内执行RPC，如果需要的话，甚至可以同时在所有这些情况下执行。
当你尝试从事务内部执行RPC时会发生什么？

如果RPC请求及其响应都要包含在事务中，并且事务的某个部分
依赖于响应返回的结果，那么就不可能使用在缓冲I/O情况下
可以使用的内存缓冲区技巧。
任何尝试采用这种缓冲方法都会使事务死锁，因为请求要到事务
保证成功后才能发送，但事务是否成功可能要到收到响应后才能
知道，如以下示例所示：

\begin{VerbatimN}[samepage=true]
begin_trans();
rpc_request();
i = rpc_response();
a[i]++;
end_trans();
\end{VerbatimN}

事务的内存足迹要在收到RPC响应之后才能确定，而在事务的内存足迹
确定之前，不可能判断事务是否能够被允许提交。
因此，与事务语义一致的唯一操作就是无条件中止事务，
这至少可以说是毫无帮助的。

以下是TM可用的一些选项：

\begin{enumerate}
\item	禁止事务内的RPC，使得任何执行RPC操作的尝试都会中止
	包含它的事务（也许还有多个嵌套事务）。
	或者，借助编译器来强制执行无RPC的事务。
	这种方法确实可行，但需要TM与其他同步原语交互。
\item	在任何给定时间只允许一个特殊的
	不可撤销事务~\cite{SpearMichaelScott2008InevitableSTM}
	进行，从而允许不可撤销事务包含RPC操作。
	这在一般情况下可行，但严重限制了RPC操作的可扩展性和性能。
	鉴于可扩展性和性能是并行性的首要目标，
	这种方法的通用性似乎有些自我限制。
	此外，使用不可撤销事务来允许RPC操作会在RPC操作开始后
	限制手动事务中止操作。
	最后，如果有一个不可撤销事务正在操作某个给定的数据项，
	那么操作同一数据项的任何其他事务都必须具有阻塞语义。
\item	识别在收到RPC响应之前就能确定事务是否成功的特殊情况，
	并在发送RPC请求之前自动将这些事务转换为不可撤销事务。
	当然，如果多个并发事务以这种方式尝试RPC调用，
	可能需要回滚除一个以外的所有事务，
	从而导致性能和可扩展性下降。
	尽管如此，对于以RPC结尾的长时间运行事务，
	这种方法可能仍然有价值。
	这种方法仍然必须限制手动事务中止操作。
\item	识别可以将RPC响应移出事务的特殊情况，
	然后使用类似于缓冲I/O的技术来处理。
\item	将事务基础设施扩展到包含RPC服务器及其客户端。
	理论上这是可能的，正如分布式数据库所展示的那样。
	然而，鉴于基于内存的TM没有慢速磁盘驱动器
	可以用来隐藏这些延迟，分布式数据库技术能否满足
	所需的性能和可扩展性要求尚不清楚。
	当然，鉴于固态硬盘的出现，数据库也很有可能需要
	重新设计其延迟隐藏方法。
\end{enumerate}

如前一节所述，I/O是TM已知的弱点，而RPC只是I/O中
一个特别棘手的情况。

\subsubsection{时间延迟}
\label{sec:future:Time Delays}

与事务外访问交互的一个重要特殊情况涉及事务内的显式时间延迟。
当然，事务内的时间延迟与TM的原子性属性相矛盾，
但这种情况可以说正是弱原子性要处理的问题。
此外，与内存映射I/O的正确交互有时需要精心控制的时序，
而应用程序经常出于各种目的使用时间延迟。
最后，可以在基于锁的临界区内执行时间延迟，
在持有hazard pointer时，在sequence-locking读侧临界区内，
以及在userspace-RCU读侧临界区内执行时间延迟，
如果需要的话，甚至可以同时在所有这些情况下执行。
从竞争或可扩展性的角度来看，这样做可能不明智，
但话说回来，这样做不会引发任何根本性的概念问题。

那么，TM如何处理事务内的时间延迟？

\begin{enumerate}
\item	忽略事务内的时间延迟。
	这看起来很优雅，但像太多其他``优雅''的解决方案一样，
	无法经受遗留代码的第一次考验。
	这样的代码可能在临界区中有重要的时间延迟，
	在被事务化时会失败。
\item	在遇到时间延迟操作时中止事务。
	这很有吸引力，但不幸的是，并不总是能够自动检测到
	时间延迟操作。
	那个紧凑的循环是在执行关键计算，还是仅仅在等待时间流逝？
\item	借助编译器来禁止事务内的时间延迟。
\item	让时间延迟正常执行。
	不幸的是，一些TM实现只在提交时才发布修改，
	这可能会违背时间延迟的目的。
\end{enumerate}

目前还不清楚是否有一个单一的正确答案。
具有弱原子性并在事务内立即发布更改（在中止时回滚这些更改）
的TM实现可能可以很好地使用最后一种方案。
即使在这种情况下，事务另一端的代码（甚至可能是硬件）
也可能需要大幅重新设计以容忍被中止的事务。
这种重新设计的需求将使得将transactional memory
应用于遗留代码变得更加困难。

\subsubsection{持久性}
\label{sec:future:Persistence}

锁原语有许多不同的类型。
一个有趣的区别是持久性，换句话说，锁是否能够独立于使用该锁的
进程的地址空间而存在。

非持久锁包括\co{pthread_mutex_lock()}、
\co{pthread_rwlock_rdlock()}以及大多数内核级锁原语。
如果实例化非持久锁数据结构的内存位置消失，锁也随之消失。
对于\co{pthread_mutex_lock()}的典型使用，这意味着当进程退出时，
其所有的锁都会消失。
这一特性可以被利用来简化程序关闭时的锁清理工作，
但使得不相关的应用程序之间共享锁变得更加困难，
因为这种共享需要应用程序共享内存。

\QuickQuiz{
	但假设一个应用程序在持有恰好位于文件映射内存区域中的
	\co{pthread_mutex_lock()}时退出了呢？
}\QuickQuizAnswer{
	确实，在这种情况下锁会持久存在，这会让其他试图获取
	这个由已不存在的进程持有的锁的进程感到沮丧。
	这就是为什么在使用位于文件映射内存区域中的\co{pthread_mutex}
	对象时需要格外小心。
}\QuickQuizEnd

持久锁有助于避免不相关应用程序之间共享内存的需求。
持久锁API包括flock系列、\co{lockf()}、System V信号量，
或\co{open()}的\co{O_CREAT}标志。
这些持久API可用于保护跨越多个应用程序运行的大规模操作，
而且在\co{O_CREAT}的情况下甚至能在操作系统重启后存活。
如果需要的话，锁甚至可以通过分布式锁管理器和分布式文件系统
跨越多个计算机系统——并在这些计算机系统中任何一个或全部重启后持续存在。

持久锁可以被任何应用程序使用，包括使用多种语言和软件环境编写的应用程序。
事实上，一个持久锁完全可能由一个用C编写的应用程序获取，
并由一个用Python编写的应用程序释放。

如何为TM提供类似的持久功能？

\begin{enumerate}
\item	将持久事务限制在专门设计来支持它们的特殊用途环境中，
	例如SQL。
	鉴于数据库系统几十年的历史，这显然是可行的，
	但不能提供持久锁所提供的同等程度的灵活性。
\item	使用某些存储设备和/或文件系统提供的快照功能。
	不幸的是，这无法处理网络通信，
	也无法处理对不提供快照功能的设备的I/O，例如U盘。
\item	制造一台时光机。
\item	完全避免这个问题，使用现有的持久设施，
	大概要避免在事务内使用这些设施。
\end{enumerate}

当然，它被称为transactional \emph{memory}这一事实
应该让我们三思，因为这个名称本身就与持久事务的概念相矛盾。
不过，将这种可能性作为探测transactional memory固有局限性的
重要测试用例来考虑仍然是值得的。

\subsection{进程修改}
\label{sec:future:Process Modification}

进程不是永恒的：
它们被创建和销毁，它们的内存映射被修改，
它们与动态库链接，它们被调试。
这些小节探讨transactional memory如何处理不断变化的执行环境。

\subsubsection{多线程事务}
\label{sec:future:Multithreaded Transactions}

在持有锁的同时创建进程和线程是完全合法的，
或者更确切地说，在持有hazard pointer时，在sequence-locking
读侧临界区内，以及在userspace-RCU读侧临界区内也是如此，
如果需要的话，甚至可以同时在所有这些情况下执行。
这不仅是合法的，而且非常简单，从以下代码片段可以看出：

\begin{VerbatimN}
pthread_mutex_lock(...);
for (i = 0; i < ncpus; i++)
	pthread_create(&tid[i], ...);
for (i = 0; i < ncpus; i++)
	pthread_join(tid[i], ...);
pthread_mutex_unlock(...);
\end{VerbatimN}

这段伪代码使用\co{pthread_create()}为每个CPU生成一个线程，
然后使用\co{pthread_join()}等待每个线程完成，
全部在\co{pthread_mutex_lock()}的保护下进行。
其效果是并行执行一个基于锁的临界区，
使用\co{fork()}和\co{wait()}也可以获得类似的效果。
当然，临界区需要足够大才能证明线程创建开销的合理性，
但在生产软件中有许多大型临界区的例子。

TM如何处理事务内的线程创建？

\begin{enumerate}
\item	声明\co{pthread_create()}在事务内是非法的，
	最好是通过中止事务来实现。
	或者，借助编译器来强制执行无\co{pthread_create()}的事务。
\item	允许在事务内执行\co{pthread_create()}，
	但只有父线程被视为事务的一部分。
	这种方法似乎与现有的和设想的TM实现合理兼容，
	但对于不知情者来说似乎是一个陷阱。
	这种方法引发了更多问题，例如如何处理冲突的子线程访问。
\item	将\co{pthread_create()}转换为函数调用。
	这种方法也是一个有吸引力的陷阱，因为它无法处理
	子线程之间相互通信这一并不罕见的情况。
	此外，它不允许事务主体的并发执行。
\item	将事务扩展为覆盖父线程和所有子线程。
	这种方法引发了关于冲突访问性质的有趣问题，
	因为父线程和子线程大概被允许相互冲突，
	但不能与其他线程冲突。
	它还引发了如果父线程在提交事务之前不等待其子线程
	会发生什么的有趣问题。
	更有趣的是，如果父线程根据参与事务的变量值
	有条件地执行\co{pthread_join()}会怎样？
	在锁的情况下，这些问题的答案相当直接。
	TM的答案留给读者作为练习。
\end{enumerate}

鉴于事务的并行执行在数据库世界中是司空见惯的，
当前的TM提案不支持这一点也许令人惊讶。
另一方面，上面的例子是锁的一种相当复杂的用法，
在简单的教科书示例中通常看不到，
所以它的缺失也许是可以预期的。
话虽如此，一些研究人员正在使用事务来自动并行化
代码~\cite{ArunRaman2010MultithreadedTransactions}，
据传其他TM研究人员正在研究事务内的fork/join并行性，
因此这个话题也许很快就会得到更全面的讨论。

\subsubsection{\tco{exec()}系统调用}
\label{sec:future:The exec System Call}

可以在基于锁的临界区内执行\co{exec()}系统调用，
在持有hazard pointer时，在sequence-locking读侧临界区内，
以及在userspace-RCU读侧临界区内执行，
如果需要的话，甚至可以同时在所有这些情况下执行。
确切的语义取决于原语的类型。

对于非持久原语（包括\co{pthread_mutex_lock()}、
\co{pthread_rwlock_rdlock()}和userspace RCU），
如果\co{exec()}成功，整个地址空间消失，
连同任何正在持有的锁。
当然，如果\co{exec()}失败，地址空间仍然存在，
因此任何相关的锁也仍然存在。
也许有点奇怪，但定义明确。

另一方面，持久原语（包括flock系列、\co{lockf()}、
System V信号量和\co{open()}的\co{O_CREAT}标志）
无论\co{exec()}成功还是失败都会继续存在，
因此被\co{exec()}执行的程序完全可以释放它们。

\QuickQuiz{
	那么由\co{mmap()}内存区域中的数据结构表示的
	非持久原语呢？
	当在这种原语的临界区内执行\co{exec()}时会发生什么？
}\QuickQuizAnswer{
	如果被\co{exec()}执行的程序映射了相同的内存区域，
	那么该程序原则上可以简单地释放锁。
	从软件工程的角度来看这种方法是否合理，
	留给读者作为练习。
}\QuickQuizEnd

当你尝试从事务内部执行\co{exec()}系统调用时会发生什么？

\begin{enumerate}
\item	禁止事务内的\co{exec()}，使得包含它的事务在遇到\co{exec()}时中止。
	这定义明确，但显然需要非TM同步原语来配合\co{exec()}使用。
\item	禁止事务内的\co{exec()}，由编译器强制执行此禁令。
	有一个C++ TM草案规范采用了这种方法，
	允许函数使用\co{transaction_safe}和\co{transaction_unsafe}
	属性进行修饰。\footnote{
		感谢Mark Moir向我指出了这个规范，
		也感谢Michael Wong此前向我指出了更早的修订版本。}
	这种方法相比在运行时中止事务有一些优势，
	但同样需要非TM同步原语来配合\co{exec()}使用。
	一个缺点是需要为大量库函数添加\co{transaction_safe}和
	\co{transaction_unsafe}属性。
\item	以类似于非持久锁原语的方式对待事务，
	使得如果\co{exec()}失败则事务继续存在，
	如果\co{exec()}成功则静默提交。
	只有部分受事务影响的变量驻留在\co{mmap()}内存中
	（因此可以在\co{exec()}系统调用成功后存活）的情况
	留给读者作为练习。
\item	如果\co{exec()}系统调用本来会成功，
	则中止事务（和\co{exec()}系统调用），
	但如果\co{exec()}系统调用本来会失败，
	则允许事务继续。
	这在某种意义上是``正确''的方法，
	但需要大量工作却得到相当不令人满意的结果。
\end{enumerate}

\co{exec()}系统调用也许是TM普遍适用性障碍中最奇特的例子，
因为目前还不完全清楚哪种方法有意义，
有些人可能会认为这仅仅是与\co{exec()}的现实交互
中的困境的反映。
话虽如此，禁止事务内\co{exec()}的两个选项
也许是这组选项中最合乎逻辑的。

类似的问题也围绕着\co{exit()}和\co{kill()}系统调用，
以及会退出事务的\co{longjmp()}或异常。
（\co{longjmp()}或异常是从哪里来的？）

\subsubsection{动态链接和加载}
\label{sec:future:Dynamic Linking and Loading}

基于锁的临界区、持有hazard pointer的代码、
sequence-locking读侧临界区和userspace-RCU读侧临界区
可以（单独或组合地）合法地包含调用动态链接和加载的函数的代码，
包括C/C++共享库和Java类库。
当然，这些库中包含的代码在编译时根据定义是不可知的。
那么，如果在事务内调用一个动态加载的函数会发生什么？

这个问题有两个部分：
\begin{enumerate*}[(a)]
\item 如何在事务内动态链接和加载一个函数，以及
\item 如何处理该函数内代码的不可知性质？
\end{enumerate*}
公平地说，问题(b)至少在理论上也给锁和userspace-RCU带来了一些挑战。
例如，动态链接的函数可能会为锁引入死锁（\IX{deadlock}），
或者可能（错误地）在userspace-RCU读侧临界区中引入
\IX{quiescent state}（静止状态）。
区别在于，虽然锁和userspace-RCU临界区中允许的操作类别已被充分理解，
但在TM的情况下似乎仍然存在相当大的不确定性。
事实上，不同的TM实现似乎有不同的限制。

那么TM如何处理动态链接和加载的库函数呢？
对于部分(a)即代码的实际加载，选项包括以下内容：

\begin{enumerate}
\item	以类似于缺页错误的方式处理动态链接和加载，
	使得函数被加载和链接，可能在此过程中中止事务。
	如果事务被中止，重试将发现函数已经存在，
	因此可以预期事务正常进行。
\item	禁止在事务内动态链接和加载函数。
\end{enumerate}

对于部分(b)即无法检测尚未加载的函数中TM不友好的操作，
可能的选项包括以下内容：

\begin{enumerate}
\item	直接执行代码：
	如果函数中有任何TM不友好的操作，则简单地中止事务。
	不幸的是，这种方法使编译器无法确定一组给定的事务
	是否可以安全地组合。
	无论如何都允许可组合性的一种方式是不可撤销事务，
	然而当前的实现在任何给定时间只允许一个不可撤销事务进行，
	这可能严重限制性能和可扩展性。
	不可撤销事务也会限制手动事务中止操作的使用。
	最后，如果有一个不可撤销事务正在操作某个给定的数据项，
	操作同一数据项的任何其他事务都不能具有非阻塞语义。
\item	修饰函数声明以指示哪些函数是TM友好的。
	这些修饰可以由编译器的类型系统来强制执行。
	当然，对于许多语言来说，这需要提出、标准化和实现语言扩展，
	带有相应的时间延迟，以及对大量原本无关的库函数
	进行相应的修饰。
	话虽如此，标准化工作已经在
	进行中~\cite{Ali-Reza-Adl-Tabatabai2009CppTM,Ali-Reza-Adl-Tabatabai2012CppTM}。
\item	如上所述，禁止在事务内动态链接和加载函数。
\end{enumerate}

I/O操作当然是TM已知的弱点，而动态链接和加载
可以被视为I/O的又一个特殊情况。
尽管如此，TM的支持者必须要么解决这个问题，
要么接受TM只是并行程序员工具箱中几个工具之一这样的现实。
（公平地说，许多TM支持者早已接受了一个不仅仅包含TM的世界。）

\subsubsection{内存映射操作}
\label{sec:future:Memory-Mapping Operations}

在基于锁的临界区内执行内存映射操作（包括\co{mmap()}、
\co{shmat()}和\co{munmap()}~\cite{TheOpenGroup1997SUS}）
是完全合法的，在持有hazard pointer时、
在sequence-locking读侧临界区内、以及在userspace-RCU读侧
临界区内也是如此，如果需要的话甚至可以同时在所有这些情况下执行。
当你尝试从事务内部执行这样的操作时会发生什么？
更关键的是，如果被重映射的内存区域包含参与当前线程事务的
某些变量会怎样？
如果这个内存区域包含参与其他线程事务的变量又会怎样？

不需要考虑TM系统的元数据被重映射的情况，
因为大多数锁原语都不定义重映射其锁变量的结果。

以下是一些TM内存映射选项：

\begin{enumerate}
\item	事务内的内存重映射是非法的，将导致所有包含它的事务被中止。
	这确实在某种程度上简化了问题，但也要求TM与那些在其临界区内
	容忍重映射的同步原语互操作。
\item	事务内的内存重映射是非法的，由编译器强制执行此禁令。
\item	事务内的内存映射是合法的，但会中止所有在被映射区域中
	拥有变量的其他事务。
\item	事务内的内存映射是合法的，但如果被映射的区域与当前事务的
	内存足迹重叠，映射操作将失败。
\item	所有内存映射操作，无论是在事务内还是事务外，
	都检查被映射区域是否与系统中所有事务的内存足迹冲突。
	如果有重叠，则内存映射操作失败。
\item	与系统中任何事务的内存足迹重叠的内存映射操作的效果
	由TM冲突管理器决定，它可以动态确定是让内存映射操作失败
	还是中止任何冲突的事务。
\end{enumerate}

值得注意的是，\co{munmap()}会使相关内存区域处于未映射状态，
这可能有额外的有趣含义。\footnote{
	映射和取消映射之间的这种差异由Josh Triplett指出。}

\subsubsection{调试}
\label{sec:future:Debugging}

通常的调试操作如断点在基于锁的临界区内和userspace-RCU读侧
临界区内正常工作。
然而，在最初的transactional-memory硬件
实现~\cite{DaveDice2009ASPLOSRockHTM}中，事务内的异常
将中止该事务，这反过来意味着断点会中止所有包含它的事务。

那么如何调试事务？

\begin{enumerate}
\item	在包含断点的事务中使用软件仿真技术。
	当然，每当在任何事务的范围内设置断点时，
	可能需要仿真所有事务。
	如果运行时系统无法确定给定的断点是否在事务的范围内，
	那么为了安全起见，可能需要仿真所有事务。
	然而，这种方法可能会带来显著的开销，
	这反过来可能会掩盖正在追踪的bug。
\item	仅使用能够处理断点异常的硬件TM实现。
	不幸的是，截至本文撰写时（2021年3月），所有此类
	实现都是研究原型。
\item	仅使用软件TM实现，它们（非常粗略地说）
	比较简单的硬件TM实现对异常更具容忍性。
	当然，软件TM的开销往往比硬件TM更高，
	因此这种方法可能并非在所有情况下都可接受。
\item	更仔细地编程，以便首先避免在事务中出现bug。
	一旦你弄清楚如何做到这一点，请务必让大家都知道这个秘诀！
\end{enumerate}

有一些理由相信transactional memory将带来与其他同步机制相比的
\IX{productivity}（生产力）改进，
但这些改进如果传统调试技术不能应用于事务的话，
似乎很容易被丧失。
如果transactional memory要被新手用于大型事务，
这一点似乎尤其正确。
相比之下，硬派的``顶尖''程序员也许能够不用这些调试辅助工具，
尤其是对于小型事务。

因此，如果transactional memory要兑现其对新手程序员的生产力承诺，
调试问题确实需要得到解决。

\subsection{同步}
\label{sec:future:Synchronization}

如果transactional memory有朝一日证明它能够满足所有人的所有需求，
那么它就不需要与任何其他同步机制交互。
在那之前，它需要与能做它不能做之事的同步机制配合工作，
或者与在给定情况下更自然地工作的同步机制配合。
以下各节概述了该领域当前的挑战。

\subsubsection{锁}
\label{sec:future:Locking}

在持有其他锁的同时获取锁是很常见的做法，这工作得很好，
至少只要使用了通常众所周知的软件工程技术来避免死锁。
从RCU读侧临界区内获取锁也不罕见，
这减轻了死锁方面的担忧，因为RCU读侧原语
不能参与基于锁的死锁循环。
在持有hazard pointer和sequence-lock读侧临界区内也可以获取锁。
但当你尝试从事务内部获取锁时会发生什么？

理论上，答案很简单：
只需将表示锁的数据结构作为事务的一部分进行操作，
一切就会完美运作。
在实践中，根据TM系统的实现细节，可能会出现许多
不明显的复杂情况~\cite{Volos2008TRANSACT}。
这些复杂情况可以解决，但代价是在事务外获取锁的开销增加45\pct，
在事务内获取锁的开销增加300\pct。
虽然这些开销对于包含少量锁的事务性程序可能是可以接受的，
但对于希望偶尔使用事务的生产级基于锁的程序来说，
这些开销往往是完全不可接受的。

以下是TM可用的一些选项：

\begin{enumerate}
\item	仅使用对锁友好的TM实现~\cite{DaveDice2006DISC}。
	不幸的是，对锁不友好的实现具有一些有吸引力的属性，
	包括成功事务的低开销和能够容纳极大事务的能力。
\item	在将TM引入基于锁的程序时仅``小规模''使用TM，
	从而适应对锁友好的TM实现的限制。
\item	完全搁置基于锁的遗留系统，用事务重新实现一切。
	这种方法不缺少倡导者，但这要求
	本系列中描述的所有问题都得到解决。
	在解决这些问题所需的时间内，竞争的同步机制
	当然也有机会得到改进。
\item	严格将TM用作基于锁系统中的优化，正如
	TxLinux~\cite{ChistopherJRossbach2007a}小组和
	大量transactional lock elision
	项目~\cite{MartinPohlack2011HTM2TLE,Kleen:2014:SEL:2566590.2576793,PascalFelber2016rwlockElision,SeongJaePark2020HTMRCUlock}所做的那样。
	这种方法似乎是合理的，但锁的设计约束
	（如需要避免死锁）仍然完全保留。
\item	努力减少施加在锁原语上的开销。
\end{enumerate}

TM和锁之间可能存在接口问题这一事实令许多人感到惊讶，
这突显了在实际生产软件中试用新机制和原语的必要性。
幸运的是，开源的出现意味着大量这样的软件现在对所有人
（包括研究人员）都是免费可用的。

\subsubsection{读写锁}
\label{sec:future:Reader-Writer Locking}

在持有其他锁的同时以读模式获取读写锁（reader-writer lock）是很常见的做法，
这直接就能工作，至少只要使用了通常众所周知的软件工程技术
来避免死锁。
从RCU读侧临界区内以读模式获取读写锁也能工作，
而且这样做减轻了死锁方面的担忧，
因为RCU读侧原语不能参与基于锁的死锁循环。
在持有hazard pointer和sequence-lock读侧临界区内也可以获取锁。
但当你尝试从事务内部以读模式获取读写锁时会发生什么？

不幸的是，在事务中以直接的方式读获取传统的基于计数器的读写锁
会违背读写锁的目的。
要理解这一点，考虑一对事务同时尝试读获取同一个读写锁。
因为读获取涉及修改读写锁的数据结构，
会产生冲突，这将回滚两个事务中的一个。
这种行为与读写锁允许并发读者的目标完全不一致。

以下是TM可用的一些选项：

\begin{enumerate}
\item	使用per-CPU或per-thread读写
	锁~\cite{WilsonCHsieh92a}，它允许给定的CPU
	（或线程）在读获取锁时仅操作本地数据。
	这将避免两个事务同时读获取锁之间的冲突，
	允许两者如预期那样继续进行。
	不幸的是，(1)~per-CPU/thread锁的写获取开销可能极高，
	(2)~per-CPU/thread锁的内存开销可能令人望而却步，
	(3)~这种转换仅在你能访问相关源代码时才可用。
	其他更新的可扩展读写锁~\cite{YossiLev2009SNZIrwlock}
	可能避免部分或全部这些问题。
\item	在将TM引入基于锁的程序时仅``小规模''使用TM，
	从而避免从事务内部读获取读写锁。
\item	完全搁置基于锁的遗留系统，用事务重新实现一切。
	这种方法不缺少倡导者，但这要求
	本系列中描述的\emph{所有}问题都得到解决。
	在解决这些问题所需的时间内，竞争的同步机制
	当然也有机会得到改进。
\item	严格将TM用作基于锁系统中的优化，正如
	TxLinux~\cite{ChistopherJRossbach2007a}小组所做的，
	以及更近期使用TM来省略读写锁的
	工作~\cite{PascalFelber2016rwlockElision}所做的那样。
	这种方法似乎是合理的，至少在\Power{8}
	CPU~\cite{Le:2015:TMS:3266491.3266500}上是这样，
	但锁的设计约束（如需要避免死锁）仍然完全保留。
\end{enumerate}

当然，将TM与读写锁结合可能还有其他不明显的问题，
就像与排他锁结合时实际上确实出现了问题一样。

\subsubsection{延迟回收}
\label{sec:future:Deferred Reclamation}

本节主要关注RCU。
将TM与其他延迟回收机制（如\IXalt{reference counters}{reference count}
和hazard pointers）结合时也会出现类似的问题和可能的解决方案。
在下面的文本中，已知的差异会被特别指出。

引用计数（reference counting）、hazard pointers和RCU都被大量使用，如
\cref{sec:defer:RCU Related Work,sec:defer:Which to Choose? (Production Use)}
中所述。
这意味着任何选择不克服本节中指出的每个挑战的TM实现
都需要与所有这些同步机制干净而高效地互操作。

德克萨斯大学奥斯汀分校的TxLinux小组似乎是
接受RCU/TM互操作性挑战的
小组~\cite{ChistopherJRossbach2007a}。
因为他们将TM应用于使用RCU的Linux 2.6内核，
所以他们别无选择，只能集成TM和RCU，用TM取代锁来进行RCU更新。
不幸的是，虽然论文确实指出RCU实现的锁
（例如\co{rcu_ctrlblk.lock}）被转换为事务，
但对于那些由基于RCU的更新使用的锁
（例如\co{dcache_lock}）做了什么却没有提及。

更近期，\ppl{Dimitrios}{Siakavaras}等人将HTM和RCU应用于搜索
树~\cite{Siakavaras2017CombiningHA,DimitriosSiakavaras2020RCU-HTM-B+Trees}，
\ppl{Christina}{Giannoula}等人使用HTM和RCU来对图进行
着色~\cite{ChristinaGiannoula2018HTM-RCU-graphcoloring}，
而\ppl{SeongJae}{Park}等人使用HTM和RCU来优化\IXacr{numa}系统上的
高竞争锁~\cite{SeongJaePark2020HTMRCUlock}。

重要的是要注意，RCU允许读者和更新者并发运行，
进一步允许RCU读者访问正在被更新的数据。
当然，RCU的这一特性，无论其性能、可扩展性和
实时响应优势如何，都与TM的底层原子性属性相矛盾，
尽管\Power{8} CPU系列的挂起事务
设施~\cite{Le:2015:TMS:3266491.3266500}使其成为这一规则的例外。

那么基于TM的更新应该如何与并发RCU读者交互？
一些可能的方式如下：

\begin{enumerate}
\item	RCU读者中止并发冲突的TM更新。
	这实际上是TxLinux项目采用的方法。
	这种方法确实保留了RCU语义，也保留了RCU的读侧性能、
	可扩展性和实时响应特性，但它有一个不幸的副作用，
	即不必要地中止冲突的更新。
	在最坏的情况下，一长串RCU读者可能饿死所有更新者，
	这在理论上可能导致系统挂起。
	此外，不是所有TM实现都提供实现此方法所需的强原子性，
	而且有充分的理由。
\item	与冲突TM更新并发运行的RCU读者从任何冲突的RCU加载中
	获取旧的（事务前的）值。
	这保留了RCU语义和性能，也防止了RCU更新饥饿（\IX{starvation}）。
	然而，不是所有TM实现都能及时提供对已被进行中事务
	暂时更新的变量旧值的访问。
	特别是，在日志中维护旧值的基于日志的TM实现
	（因此提供了出色的TM提交性能）不太可能喜欢这种方法。
	也许可以利用\co{rcu_dereference()}原语来允许RCU
	在更广泛的TM实现中访问旧值，
	尽管性能可能仍然是个问题。
	尽管如此，有流行的TM实现已经以这种方式与RCU
	集成~\cite{DonaldEPorter2007TRANSACT,PhilHoward2011RCUTMRBTree,
	PhilipWHoward2013RCUrbtree}。
\item	如果RCU读者执行了与进行中事务冲突的访问，
	那么该RCU访问将被延迟，直到冲突事务提交或中止。
	这种方法保留了RCU语义，但不保留RCU的性能或实时响应，
	特别是在存在长时间运行事务的情况下。
	此外，不是所有TM实现都能够延迟冲突的访问。
	尽管如此，对于仅支持小型事务的硬件TM实现，
	这种方法似乎非常合理。
\item	RCU读者被转换为事务。
	这种方法基本上保证了RCU与任何TM实现兼容，
	但它也将TM的回滚强加于RCU读侧临界区，
	破坏了RCU的实时响应保证，也降低了RCU的读侧性能。
	此外，在任何RCU读侧临界区包含相关TM实现无法处理的
	操作的情况下，这种方法是不可行的。
	这种方法更难应用于hazard pointers和引用计数，
	因为它们没有将读者作为代码段的明确定义。
\item	RCU的许多更新侧用法修改单个指针来发布新的数据结构。
	在其中一些情况下，只要事务尊重内存排序并且回滚过程
	使用\co{call_rcu()}来释放相应的结构，
	就可以安全地允许RCU看到随后被回滚的事务性指针更新。
	不幸的是，不是所有TM实现都尊重事务内的内存屏障。
	显然，人们的想法是，因为事务应该是原子的，
	事务内访问的顺序不应该有影响。
\item	禁止在RCU更新中使用TM。
	这保证可以工作，但限制了TM的使用。
\end{enumerate}

可能还会发现更多的方法，特别是考虑到用户级RCU和
hazard-pointer实现的出现。\footnote{
	感谢TxLinux小组、Maged Michael和Josh Triplett
	提出了上述多种替代方案。}
有趣的是，许多性能和可扩展性较好的STM实现在内部使用了
类RCU技术~\cite{UCAM-CL-TR-579,KeirFraser2007withoutLocks,Gu:2019:PSE:3358807.3358885,Kim:2019:MSR:3297858.3304040}。

\QuickQuiz{
	MV-RLU看起来相当不错！
	它是不是完胜RCU？
}\QuickQuizAnswer{
	快速阅读摘要可能会给人这种印象，
	但更仔细的读者会注意到最后一句中的``对于广泛的工作负载''
	这一短语。
	事实证明这个短语相当重要：

	\begin{enumerate}
	\item	他们的RCU评估使用了同步宽限期，这不必要地
		限制了更新速度，正如他们的第6.2.1节所述。
		参见本书的\cref{fig:datastruct:Read-Side RCU-Protected Hash-Table Performance For Schroedinger's Zoo in the Presence of Updates}
		\cpageref{fig:datastruct:Read-Side RCU-Protected Hash-Table Performance For Schroedinger's Zoo in the Presence of Updates}，
		可以看到历史悠久的异步\co{call_rcu()}原语
		使RCU能够在大量更新者的情况下表现和扩展得相当好。
		此外，在他们论文的第3.7节中，作者承认
		异步宽限期对MV-RLU的可扩展性很重要。
		公平的比较也应该允许RCU享受异步的好处。
	\item	他们使用了一个调优不佳的含10,000个元素的
		1,000桶哈希表。
		此外，他们的448个硬件线程需要远多于1,000个桶
		才能避免他们正确指出的限制RCU在其基准测试中
		性能的锁竞争（lock contention）。
		有用的比较应该使用适当调优的哈希表。
	\item	他们的RCU哈希表使用了per-bucket锁，
		他们将其称为瓶颈，鉴于长哈希链和
		桶与线程的小比率，这并不令人惊讶。
		他们的一些竞争机制改为使用无锁（lock-free）技术，
		从而避免了per-bucket锁的瓶颈，
		愤世嫉俗者可能会声称这揭示了作者在其他方面
		令人费解地选择调优不佳哈希表的原因。
		作者图4中间行的第一张图显示了RCU在不受
		人为瓶颈束缚时能达到的效果，
		同一行第二张图的前半部分也是如此。
	\item	他们的链表操作允许RLU对列表中不同元素进行并发修改，
		而RCU被迫串行化更新。
		同样，RCU一直与无锁更新者配合得很好，
		这一事实已在作者引用的学术文献中
		阐述~\cite{MathieuDesnoyers2012URCU}。
		公平的比较应该对RCU使用与MV-RLU相同风格的更新。
	\item	作者没有考虑将RCU和sequence locking结合使用，
		Linux内核中使用这种组合来给读者
		提供多指针更新的一致视图。
	\item	作者没有考虑基于RCU的Issaquah
		Challenge~\cite{PaulEMcKenney2016IssaquahCPPCON}解决方案，
		该方案也给读者提供多指针更新的一致视图，
		尽管对``一致''的定义较弱。
	\end{enumerate}

	令人惊讶的是，这篇论文的匿名审稿人没有要求
	对MV-RLU和RCU进行同类比较。
	尽管如此，应该祝贺作者撰写了一篇学术论文，
	其中展示了良好可扩展性与强读侧一致性相结合的罕见例子。
	也应该祝贺他们克服了传统学术界对异步宽限期的偏见，
	这极大地帮助了他们的可扩展性。

	有趣的是，RLU和RCU采用不同的方法来避免
	\ppl{Hagit}{Attiya}等人指出的STM固有
	限制~\cite{Attiya:2009:STMReadOnlyLimits}。
	RCU避免提供严格的可串行化，RLU避免提供
	不可见的只读事务，两者因此都避免了这些限制。
}\QuickQuizEnd

\subsubsection{事务外访问}
\label{sec:future:Extra-Transactional Accesses}

在基于锁的临界区内，操作在该锁的临界区外被并发访问
甚至修改的变量是完全合法的，一个常见的例子是统计计数器（statistical counter）。
在RCU读侧临界区内也可以做同样的事情，
而且这实际上是常见情况。

鉴于生产数据库系统中普遍存在的所谓``脏读''等机制，
TM的支持者认真关注事务外访问就不足为奇了，
弱原子性~\cite{Blundell2006TMdeadlock}的概念就是一个例子。

以下是一些事务外访问选项：

\begin{enumerate}
\item	由事务外访问引起的冲突总是中止事务。
	这就是强原子性。
\item	由事务外访问引起的冲突被忽略，
	因此只有事务之间的冲突才能中止事务。
	这就是弱原子性。
\item	允许事务在特殊情况下执行非事务操作，
	例如在分配内存或与基于锁的临界区交互时。
\item	产生允许某些操作（例如加法）由多个事务
	在单个变量上并发执行的硬件扩展。
\item	向transactional memory引入弱语义。
	一种方法是\cref{sec:future:Deferred Reclamation}中描述的
	与RCU的结合，而Gramoli和Guerraoui
	调查了许多其他弱事务
	方法~\cite{Gramoli:2014:DTP:2541883.2541900}，例如，
	将大型``弹性''事务限制性地分割为较小的事务，
	从而降低冲突概率（尽管性能和可扩展性不温不火）。
	也许进一步的经验将表明事务外访问的某些用途
	可以被弱事务所取代。
\end{enumerate}

看起来事务是在真空中构想出来的，
不需要与任何其他同步机制交互。
如果是这样，当事务与非事务访问结合时出现大量混乱和复杂性
就不足为奇了。
但除非事务被限制在对孤立数据结构的小型更新上，
或者被限制在不与大量现有并行代码交互的新程序上，
否则如果事务要在近期产生大规模实际影响，
就绝对必须进行这种结合。

\subsection{其他事务}
\label{sec:future:Other Transactions}
% If we ever get epigraphs down at this level, perhaps the old
% "Hell is other people" quote by Jean-Paul Sartre.  His clarification
% also applies, which can be stated as "Heaven is also other people".

因为事务对彼此应该表现为原子的，
当多个事务同时尝试访问同一个变量时必须采取一些措施。
如果所有事务都在读取该变量而没有一个在更新它，
那么可以``什么都不做''，否则事务必须至少部分串行化。
这种串行化通常通过中止和回滚那些事务的某个子集来实现。

如何选择中止和回滚哪些事务？

关于这个问题已经有大量的工作。
答案是``使用竞争管理器''，但这只是将问题转移到了
``竞争管理器应该做什么？''
对竞争管理器的全面论述超出了本节的范围。
因此本节将限于该领域的几篇开创性论文。

Herlihy等人~\cite{HerlihyLMS03}
让写事务无条件地中止读事务，
这有简单性的优点，但可能导致读者饥饿。
这篇论文还引入了早期释放的概念，
它可以允许一次长遍历与小更新并发进行。
从这个角度来看，可以将RCU读者描述为``即时释放''，
但这是一种使RCU读者以完全不同于STM方式行为的特性。

Hammond等人~\cite{LanceHammond2004a}
建议使用小事务，这确实降低了竞争管理器的复杂性，
但也未能兑现STM自动并发的全部承诺。
这篇论文还引入了注解来标识不同类型的事务。

Scherer和Scott~\cite{WilliamNSchererIII2005}
描述了许多STM竞争管理策略。
这篇论文为事务外访问定义了``可见''和``不可见''的读者，
并讨论了许多竞争管理策略，包括：

\begin{enumerate}
\item	指数退避，他们得出结论认为需要手动调优。
\item	优先中止完成工作量最少的事务，
	其中之前被中止的尝试完成事务时所做的工作
	会计入该事务下次重试的工作量。
	这有很好的公平性属性，但可能导致大型事务
	长时间占用系统。
\item	优先中止完成工作量最少的事务，
	但也将等待给定事务完成的任何其他事务所做的工作
	计入该事务。
	这在组合中加入了一种优先级继承的形式。
\item	跟踪给定事务被另一个事务中止的次数，
	作为促进某种公平性概念的另一种方式。
\item	使用事务开始时间戳，以便优先中止最近开始的事务。
\item	使用事务开始时间戳，但也添加给定事务上次运行时间的指示，
	以防止被抢占的事务中止在更高优先级进程中运行的其他事务。
\end{enumerate}

这篇2005年的论文因此早期指出了竞争管理中固有的复杂性。

Porter和Witchel~\cite{DonaldEPorter2007TRANSACT,Ramadan:2008:DTM:1521747.1521799}
考虑放宽STM一致性要求以简化竞争管理。
他们指出RCU是更新会导致并发读取访问数据旧版本的一种情况。

竞争管理在HTM中也是一个问题，将在
\cref{sec:future:Conflict Handling,sec:future:Aborts and Rollbacks}
中讨论。

\QuickQuiz{
	为什么不与时俱进，将机器学习应用于竞争管理？
}\QuickQuizAnswer{
	许多事务的开销在亚微秒级别，因此在机器学习上消耗
	比通过更高效地执行事务所节省的更多CPU是太容易的事情。
	尽管如此，机器学习可能在将竞争管理器调优到
	特定工作负载方面发挥作用。
	当然，不需要调优会更好，但有时``更好''是不可实现的。
}\QuickQuizEnd

\subsection{案例研究：
	    Sequence Locking}
\label{sec:future:Case Study: Sequence Locking}

\cref{sec:defer:Sequence Locks}中描述的Sequence locking
有时被认为是STM的一种简单形式。
Sequence-locking读侧临界区通常限于执行加载操作，
每当与sequence-locking更新者发生冲突时就会被重试。
至关重要的是，sequence-locking读侧临界区之间不会相互冲突。
因此，考虑sequence locking如何处理前面各节中指出的
STM挑战是很有启发性的。

Sequence locking通过简单地在每次通过其读侧临界区时
执行TM不友好的操作来处理与外部世界的交互
（\cref{sec:future:Outside World}）。
例如，如果执行I/O写入的临界区被重试了三次，
它将执行四次该I/O写入，一次是原始通过临界区的执行，
三次是每次重试各一次。

在sequence-locking读侧临界区内执行I/O操作
可能看起来有些非常规。
然而，出于历史记录、调试或审计目的
向日志进行I/O写入（\cref{sec:future:I/O Operations}）
是非常有意义的。
RPC操作（\cref{sec:future:RPC Operations}）可以收集所需数据。
时间延迟（\cref{sec:future:Time Delays}）可以为调试目的
提供有价值的服务，例如测试重试代码路径。
即使是持久操作（\cref{sec:future:Persistence}）也可能有用，
例如执行跨应用程序的互斥。
或者，就像STM一样，sequence locking的用户可以选择
缓冲I/O并在读侧临界区结束后立即执行它。

Sequence locking以相同的方式处理进程修改
（\cref{sec:future:Process Modification}），
即在其读侧临界区的每次执行中简单地执行进程修改操作。

在临界区内创建进程和线程
（\cref{sec:future:Multithreaded Transactions}）
正常运作，尽管应该注意防止经常重试的临界区成为``fork炸弹''。
\co{exec()}系统调用（\cref{sec:future:The exec System Call}）
有完全合理的（虽然相当激进的）语义。
对动态库中函数的调用
（\cref{sec:future:Dynamic Linking and Loading}），
例如Linux内核可加载模块，与对内置函数的调用一样正常工作。
当然，如果这样的调用导致模块加载，
这可能导致临界区第一次执行时的高重试概率。
重映射内存（\cref{sec:future:Memory-Mapping Operations}）
行为直观，至少假设组成临界区的代码和sequence lock本身
都没有被重映射。
请注意，这些相同的限制也适用于普通锁。
调试器（\cref{sec:future:Debugging}）
在sequence-locking读侧临界区内正常工作，
就像在普通锁的临界区内一样。

同步原语（\cref{sec:future:Synchronization}）
可以在sequence-locking读侧临界区内自由使用。
锁（\cref{sec:future:Locking}）、
读写锁（\cref{sec:future:Reader-Writer Locking}）、
延迟回收（\cref{sec:future:Deferred Reclamation}）
以及事务外访问
（\cref{sec:future:Extra-Transactional Accesses}）
都可以在sequence-locking读侧临界区内使用，
并且都给出预期的结果。

Sequence locking使用一个简单的竞争管理器
（\cref{sec:future:Other Transactions}）。
Sequence-locking读者在有并发更新者的情况下总是被回滚
（正如Herlihy等人~\cite{HerlihyLMS03}所做的），
sequence-locking更新者通常将其竞争管理委托给某种形式的互斥，
而sequence-locking读者从不冲突，
因此从不需要管理其不存在的竞争。
这种粗略直接的方法意味着sequence-locking更新者可能饿死读者，
因此sequence locking的用户必须注意避免频繁更新。

简而言之，sequence locking相当好地处理了这些STM挑战。
这可能是sequence locking在生产中被大量使用的原因之一。
然而，sequence locking的一个关键使能因素是
将许多这些挑战推回给用户。
这种方法是合理的，因为sequence locking没有成为
统治所有同步机制的唯一机制的野心。
因此，也许广泛采用STM的最大障碍是
其支持者对完全通用性的过度渴望。

\subsection{讨论}
\label{sec:future:Discussion}

TM普遍采用的障碍导致以下结论：

\begin{enumerate}
\item	TM的一个有趣属性是事务可能被回滚和重试。
	这一属性是TM在不可逆操作方面困难的根源，
	包括非缓冲I/O、RPC、内存映射操作、
	时间延迟和\co{exec()}系统调用。
	这一属性还有一个不幸的后果，即引入了失败可能性
	中固有的所有复杂性，而且通常以开发者可见的方式体现。
\item	TM的另一个有趣属性，由Shpeisman
	等人~\cite{TatianaShpeisman2009CppTM}指出，是TM将同步
	与它保护的数据交织在一起。
	这一属性是TM在I/O、内存映射操作、事务外访问
	和调试断点方面问题的根源。
	相比之下，传统的同步原语，包括锁和RCU，
	在同步原语和它们保护的数据之间保持清晰的分离。
\item	TM领域许多工作者的既定目标之一是
	简化大型顺序程序的并行化。
	因此，单个事务通常被期望串行执行，
	这可能在很大程度上解释了TM在多线程事务方面的问题。
\end{enumerate}

\QuickQuiz{
	鉴于像\co{spin_trylock()}这样的东西，
	声称TM引入了失败的概念怎么可能有道理？？？
}\QuickQuizAnswer{
	使用锁时，\co{spin_trylock()}是一种选择，
	相应的无失败选择是\co{spin_lock()}，
	它在通常情况下使用，即在v5.11 Linux内核中
	对\co{spin_lock()}的调用是对\co{spin_trylock()}调用的
	100多倍。
	使用TM时，唯一的无失败选择是不可撤销事务，
	它在通常情况下不被使用。
	事实上，不可撤销事务甚至并非在所有TM实现中都可用。
}\QuickQuizEnd

TM研究人员和开发者应该如何应对所有这些问题？

一种方法是关注小规模的TM，聚焦于硬件辅助可能比其他同步原语
提供显著优势的小型事务，以及有证据表明TM与锁结合的方法能够
提高生产力的小型程序~\cite{VPankratius2011TMvsLockingProductivity}。
Sun在其Rock研究CPU~\cite{DaveDice2009ASPLOSRockHTM}中
采用了小事务方法。
一些TM研究人员似乎同意这两种``小即是美''的
方法~\cite{JMStone93}，其他人对TM寄予更高期望，
而还有一些人暗示高TM抱负可能是TM最大的
敌人~\cite[Section 6]{Attiya:2010:ICT:1835698.1835699}。
尽管如此，TM完全有可能能够处理更大的问题，
而本节列出了如果TM要实现这一崇高目标
必须解决的几个问题。

当然，所有参与者都应该将此视为学习经历。
看起来TM研究人员有很多东西需要向那些
使用传统同步原语成功构建了大型软件系统的
实践者学习。

反之亦然。

\QuickQuiz{
	有什么要学的？
	为什么不直接对基于内存的数据结构使用TM，
	对那些涉及本节列出的许多愚蠢极端情况的
	罕见情况使用锁？
}\QuickQuizAnswer{
	2005年刚打来电话，它想要回它那白热化的TM营销炒作。

	在2023年，TM仍然有很多需要证明的地方，
	即使有了HTM的出现，这将在接下来的
	\cref{sec:future:Hardware Transactional Memory}中介绍。
}\QuickQuizEnd

\begin{figure}
\centering
\resizebox{3in}{!}{\includegraphics{cartoons/TM-the-vision}}
\caption{STM的愿景}
\ContributedBy{fig:future:The STM Vision}{Melissa Broussard}
\end{figure}

\begin{figure}
\centering
\resizebox{2.7in}{!}{\includegraphics{cartoons/TM-the-reality-conflict}}
\caption{STM的现实：冲突}
\ContributedBy{fig:future:The STM Reality: Conflicts}{Melissa Broussard}
\end{figure}

\begin{figure}
\centering
\resizebox{3in}{!}{\includegraphics{cartoons/TM-the-reality-nonidempotent}}
\caption{STM的现实：不可撤销操作}
\ContributedBy{fig:future:The STM Reality: Irrevocable Operations}{Melissa Broussard}
\end{figure}

\begin{figure}
\centering
\resizebox{2.7in}{!}{\includegraphics{cartoons/TM-the-reality-realtime}}
\caption{STM的现实：实时响应}
\ContributedBy{fig:future:The STM Reality: Realtime Response}{Melissa Broussard}
\end{figure}

但就目前而言，STM的当前状态最好用一系列漫画来总结。
首先，\cref{fig:future:The STM Vision}展示了STM的愿景。
一如既往，现实要更加微妙，如
\cref{fig:future:The STM Reality: Conflicts,%
fig:future:The STM Reality: Irrevocable Operations,%
fig:future:The STM Reality: Realtime Response}
所异想天开地描绘的那样。\footnote{
	最近的学术研究工作正在调查用于实时用途的基于锁的STM
	系统~\cite{JimAnderson2019STMRT,CatherineNemitz2018LockSTMrealtime}，
	尽管没有任何性能结果，并且有一些迹象表明
	实时混合STM/HTM系统必须在快速的常见情况性能和
	最坏情况前进保证之间做出
	选择~\cite{DBLP:journals/corr/AlistarhKKRS14,MartinSchoeberl2010realtimeTM}。}
不那么异想天开的STM回顾也是
可用的~\cite{JoeDuffy2010RetroTM,JoeDuffy2010RetroTM2}。

或者，有人可能会认为sequence locking构成了STM的一种
在实践中被大量使用的受限形式，如
\cref{sec:future:Case Study: Sequence Locking}中所讨论的。

无论sequence locking是否被认为是STM在实际实践中的旗手，
一些商用硬件支持HTM的受限变体，
这将在下一节中讨论。

\IfTwoColumn{}{\FloatBarrier}
% future/htm.tex
% mainfile: ../perfbook.tex
% SPDX-License-Identifier: CC-BY-SA-3.0

\section{硬件事务内存}
\label{sec:future:Hardware Transactional Memory}
%
\epigraph{在自动化之前，确保你的报告系统合理地清晰和高效。
	  否则，你的新电脑只会加速这团混乱。}
	 {Robert Townsend}
% If at first you do succeed---try to hide your astonishment.
% Harry F.~Banks

截至2021年，\IXacrf{htm}已经在多种商用通用计算机系统上可用多年~\cite{Yoo:2013:PEI:2503210.2503232,RickMerrit2011PowerTM,ChristianJacobi2012MainframeTM,TimothyHayes2020ARM-HTM}。
本节试图确定HTM在并行程序员工具箱中的位置。

从概念上看，HTM利用处理器缓存和推测执行，
使一组指定的语句（即"事务"）从其他处理器上运行的
任何其他事务的角度来看是原子生效的。
该事务由begin-transaction机器指令启动，
并由commit-transaction机器指令完成。
通常还有一条abort-transaction机器指令，
它会取消推测（如同begin-transaction指令及其后续所有指令都未执行过一样），
并从失败处理程序处开始执行。
失败处理程序的位置通常由begin-transaction指令指定，
可以是显式的失败处理程序地址，也可以通过该指令本身设置的条件码来指定。
每个事务相对于所有其他事务原子执行。

HTM具有多项重要优势，包括自动动态划分数据结构、
减少同步原语的缓存未命中（cache miss），以及支持相当数量的实际应用。

然而，阅读细则总是值得的，HTM也不例外。
本节的一个重点是确定在什么条件下HTM的优势超过其细则中隐藏的复杂性。
为此，\cref{sec:future:HTM Benefits WRT Locking}
描述了HTM的优势，
\cref{sec:future:HTM Weaknesses WRT Locking}描述了其弱点。
这与早期论文~\cite{McKenney2007PLOSTM,PaulEMcKenney2010OSRGrassGreener}
以及前一节中使用的方法相同。\footnote{
	我衷心感谢与其他作者Maged Michael、Josh Triplett
	和Jonathan Walpole以及Andi Kleen进行的许多富有启发的讨论。}

\Cref{sec:future:HTM Weaknesses WRT Locking When Augmented}接着描述了
HTM相对于Linux内核（以及许多用户空间应用程序）中使用的同步原语组合的弱点。
\Cref{sec:future:Where Does HTM Best Fit In?}探讨了HTM
最适合在并行程序员工具箱中的位置，
\cref{sec:future:Potential Game Changers}列出了一些可能
大大增加HTM适用范围和吸引力的事件。
最后，\cref{sec:future:Conclusions}
给出了总结性评述。

\subsection{HTM 相对于锁的优势}
\label{sec:future:HTM Benefits WRT Locking}

HTM的主要优势是：
（1）避免了其他同步原语经常导致的缓存未命中，
（2）能够动态划分数据结构，
以及（3）它具有相当数量的实际应用。
我打破TM的传统，没有单独列出易用性，原因有二。
首先，易用性应该源于HTM的主要优势，这正是本节的重点。
其次，围绕着测试原始编程天赋的尝试存在相当大的争议~\cite{RichardBornat2006SheepGoats,SaeedDehnadi2009SheepGoats,ElizabethPatitsas2020GradesNotBimodal}，
甚至在工作面试中使用小型编程练习也存在争议~\cite{RegBraithwaite2007FizzBuzz}。
这表明我们确实没有牢固地掌握是什么使编程变得容易或困难。
因此，本节的其余部分重点讨论上面列出的三个优势。

\subsubsection{避免同步 Cache Miss}
\label{sec:future:Avoiding Synchronization Cache Misses}

大多数同步机制基于由原子指令操作的数据结构。
因为这些原子指令通常首先需要使相关的\IX{cache line}被运行它们的CPU所拥有，
随后在其他CPU上执行同一同步原语实例将导致缓存未命中。
这些通信缓存未命中严重降低了传统同步机制的性能和可扩展性~\cite[Section 4.2.3]{Anderson97}。

相比之下，HTM通过使用CPU的缓存进行同步，避免了对单独的同步数据结构
及其导致的缓存未命中的需要。
HTM的优势在锁数据结构被放置在单独的缓存行中的情况下最大，
在这种情况下，将给定的临界区转换为HTM事务可以减少一次完整的缓存未命中开销。
这些节省对于短临界区的常见情况可能非常显著，至少在被省略的锁
没有与该锁保护的频繁写入变量共享缓存行的情况下如此。

\QuickQuiz{
	为什么频繁写入的变量与锁变量共享缓存行会有影响？
}\QuickQuizAnswer{
	如果锁与它保护的某些变量在同一缓存行中，
	那么一个CPU对这些变量的写入将使所有其他CPU的该缓存行失效。
	这些\IXpl{invalidation}将产生大量冲突和重试，甚至可能
	降低性能和可扩展性，比使用锁还差。
}\QuickQuizEnd

\subsubsection{数据结构的动态分区}
\label{sec:future:Dynamic Partitioning of Data Structures}

使用某些传统同步机制的一个主要障碍是需要静态划分数据结构。
有许多数据结构可以轻松划分，最突出的例子是哈希表，
其中每个哈希链构成一个分区。
为每个哈希链分配一个锁，就可以轻松地将哈希表中局限于给定链的操作并行化。\footnote{
	同样容易将此方案扩展到访问多个哈希链的操作，
	只需让这些操作按哈希顺序获取所有相关链的锁即可。}
对于数组、基数树、跳表和其他一些数据结构，划分同样很简单。

然而，对于许多类型的树和图，划分相当困难，结果通常也相当复杂~\cite{Ellis80}。
虽然可以使用两阶段锁和哈希锁数组来划分一般数据结构，
但其他技术已被证明更可取~\cite{DavidSMiller2006HashedLocking}，
如将在\cref{sec:future:HTM Weaknesses WRT Locking When Augmented}中讨论的。
鉴于其避免了同步缓存未命中的特性，
HTM因此对于大型不可划分的数据结构是一个非常现实的选择，
至少在更新相对较小的情况下是如此。

\QuickQuiz{
	为什么相对较小的更新对\IXacr{htm}的性能和可扩展性很重要？
}\QuickQuizAnswer{
	更新越大，冲突的概率就越大，因此重试的概率也越大，
	这会降低性能。
}\QuickQuizEnd

\subsubsection{实际价值}
\label{sec:future:Practical Value}

HTM的实际价值已在多个硬件平台上得到了一定程度的证明，包括
Sun Rock~\cite{DaveDice2009ASPLOSRockHTM}、
Azul Vega~\cite{CliffClick2009AzulHTM}、
IBM Blue Gene/Q~\cite{RickMerrit2011PowerTM}、
Intel Haswell TSX~\cite{RaviRajwar2012TSX}和
IBM System z~\cite{ChristianJacobi2012MainframeTM}。

预期的实际优势包括：

\begin{enumerate}
\item	内存数据访问和更新的锁省略~\cite{Martinez01a,Rajwar02a}。
\item	对大型不可划分数据结构的并发访问和小规模随机更新。
\end{enumerate}

然而，HTM也有一些非常实际的不足，将在下一节中讨论。

\subsection{HTM 相对于锁的劣势}
\label{sec:future:HTM Weaknesses WRT Locking}

HTM的概念非常简单：
一组对内存的访问和更新原子地发生。
然而，与许多简单的想法一样，当你将其应用于现实世界中的真实系统时，
就会出现复杂情况。
这些复杂情况如下：

\begin{enumerate}
\item	事务大小限制。
\item	冲突处理。
\item	中止和回滚。
\item	缺乏前进进度保证。
\item	不可撤销操作。
\item	语义差异。
\end{enumerate}

以下各节将分别讨论这些复杂情况，然后给出总结。

\subsubsection{事务大小限制}
\label{sec:future:Transaction-Size Limitations}

当前HTM实现的事务大小限制源于使用处理器缓存来保存
受事务影响的数据。
虽然这允许给定的CPU通过在其缓存范围内执行事务来使事务对其他CPU
显得原子，但这也意味着任何不适合的事务都无法提交。
此外，改变执行上下文的事件，如中断、系统调用、异常、陷阱和上下文切换，
要么必须中止相关CPU上正在进行的事务，要么必须由于其他执行上下文的
缓存占用而进一步限制事务大小。

当然，现代CPU往往具有大型缓存，许多事务所需的数据可以轻松放入
一兆字节的缓存中。
不幸的是，对于缓存来说，单纯的大小并不是唯一重要的因素。
问题在于大多数缓存可以被认为是用硬件实现的哈希表。
然而，硬件缓存不会链接其桶（通常称为\emph{组}），而是
为每组提供固定数量的缓存行。
给定缓存中每组提供的元素数量
称为该缓存的\emph{\IXalt{关联度}{cache associativity}}。

虽然缓存关联度各不相同，但我正在打字的这台笔记本电脑上
L0缓存的八路关联并不罕见。
这意味着如果一个给定的事务需要访问九个缓存行，
且所有九个缓存行都映射到同一组，
那么该事务不可能完成，不管缓存中还有多少兆字节的额外空间。
是的，对于给定数据结构中随机选择的数据元素，
该事务能够提交的概率相当高，
但不能保证~\cite{PaulEMcKenney2012HTMCacheGeometry}。

已有一些研究工作试图缓解这一限制。
全关联\emph{victim缓存}可以缓解关联度约束，
但目前对victim缓存的大小有严格的性能和能效约束。
也就是说，用于未修改缓存行的HTM victim缓存可以相当小，
因为它们只需要保留地址：
数据本身可以写回内存或被其他缓存影子保存，
而地址本身就足以检测冲突写入~\cite{RaviRajwar2012TSX}。

\IXAcrmfst{utm}
方案~\cite{CScottAnanian2006,KevinEMoore2006}
使用DRAM作为极大的victim缓存，但将这些方案集成到
生产级\IXalt{cache-coherence}{cache coherence}机制中
仍然是一个未解决的问题。
此外，使用DRAM作为victim缓存可能会带来不利的
性能和能效后果，特别是如果victim缓存要求
\IXalth{全关联}{fully associative}{cache}的话。
最后，UTM的"无界"方面假设所有DRAM
都可以用作victim缓存，而实际上分配给给定CPU的
大量但仍然固定的DRAM将限制该CPU事务的大小。
其他方案使用硬件和软件transactional memory的组合~\cite{SanjeevKumar2006}，
人们可以想象使用\IXacr{stm}作为HTM的后备机制。

然而，据我所知，除了缩略表示TM读集之外，
当前可用的系统没有实现这些研究想法中的任何一个，
也许是有充分理由的。

\subsubsection{冲突处理}
\label{sec:future:Conflict Handling}

第一个复杂因素是\emph{冲突}的可能性。
例如，假设事务~A和~B定义如下：

\begin{VerbatimU}
Transaction A       Transaction B

x = 1;              y = 2;
y = 3;              x = 4;
\end{VerbatimU}

假设每个事务在各自的处理器上并发执行。
如果事务~A在事务~B存储\co{y}的同时存储\co{x}，
则两个事务都无法继续。
为了理解这一点，假设事务~A执行对\co{y}的存储。
那么事务~A将交错在事务~B内部，违反了事务之间
原子执行的要求。
允许事务~B执行对\co{x}的存储同样违反原子执行要求。
这种情况被称为\emph{冲突}，当两个并发事务访问同一变量
且至少有一个访问是存储时就会发生。
因此系统有义务中止一个或两个事务以允许执行继续。
究竟中止哪个事务是一个有趣的课题，
很可能在未来相当长时间内仍能产出博士论文，
例如参见~\cite{EgeAkpinar2011HTM2TLE}。\footnote{
	Liu和Spear题为"Toxic
	Transactions"~\cite{YujieLiu2011ToxicTransactions}的论文
	特别具有教育意义。}
就本节而言，我们可以假设系统做出随机选择。

另一个复杂因素是冲突检测，至少在最简单的情况下比较直接。
当处理器正在执行事务时，它会标记该事务触及的每个缓存行。
如果处理器的缓存收到涉及已被当前事务标记为已触及的缓存行的请求，
则发生了潜在冲突。
更复杂的系统可能会尝试将当前处理器的事务排序在发送请求的处理器的事务之前，
优化此过程也可能在相当长时间内产出博士论文。
然而本节假设一种非常简单的冲突检测策略。

然而，要使HTM有效工作，冲突概率必须相当低，这反过来要求数据结构的组织方式
能维持足够低的冲突概率。
例如，具有简单插入、删除和搜索操作的红黑树符合此描述，
但维护准确元素计数的红黑树则不符合。\footnote{
	更新计数的需要将导致树的添加和删除操作相互冲突，
	产生强不可交换性~\cite{HagitAttiya2011LawsOfOrder,Attiya:2011:LOE:1925844.1926442,PaulEMcKenney2011SNC}。}
另一个例子是，在单个事务中枚举树中所有元素的红黑树将具有很高的冲突概率，
从而降低性能和可扩展性。
因此，许多串行程序在HTM能有效工作之前需要一些重构。
在某些情况下，实践者会倾向于采取额外步骤
（在红黑树的情况下，也许切换到可划分的数据结构如基数树或哈希表），
直接使用锁，特别是在HTM在所有相关架构上广泛可用之前~\cite{CliffClick2009AzulHTM}。

\QuickQuiz{
	无论选择什么同步机制，红黑树怎么可能高效地枚举树中的所有元素？
}\QuickQuizAnswer{
	在许多情况下，枚举不需要精确。
	在这些情况下，可以使用hazard pointer或\IXacr{rcu}来保护读者，
	这提供了与任何给定插入或删除操作的低冲突概率。
}\QuickQuizEnd

此外，并发事务之间潜在的冲突访问可能导致失败。
处理此类失败将在下一节中讨论。

\subsubsection{中止和回滚}
\label{sec:future:Aborts and Rollbacks}

因为任何事务都可能在任何时候被中止，
所以事务不能包含无法回滚的语句，这一点很重要。
这意味着事务不能执行I/O、系统调用或调试断点
（HTM事务中不能在调试器中单步执行！！！）。
相反，事务必须将自己限制在访问正常的缓存内存。
此外，在某些系统上，中断、异常、陷阱、TLB未命中和其他事件也会中止事务。
鉴于因不当处理错误条件而导致的大量错误，
有理由问中止和回滚对易用性有什么影响。

\QuickQuiz{
	但为什么调试器不能通过在事务的连续行上设置断点来模拟单步执行，
	依赖重试来重新追踪事务早期实例的步骤？
}\QuickQuizAnswer{
	这种方案可能以相当高的概率工作，但它
	可能以令大多数用户相当惊讶的方式失败。
	为了理解这一点，考虑以下事务：

\begin{fcvlabel}[ln:future:htm:debug rollbacks]
\begin{VerbatimN}[commandchars=\\\[\]]
begin_trans();
if (a) {
	do_one_thing();
	do_another_thing();	\lnlbl[another]
} else {
	do_a_third_thing();
	do_a_fourth_thing();
}
end_trans();
\end{VerbatimN}
\end{fcvlabel}

	\begin{fcvref}[ln:future:htm:debug rollbacks]
	假设用户在\clnref{another}处设置了断点，
	该断点触发，中止事务并进入调试器。
	\end{fcvref}
	假设在断点触发和调试器停止所有线程之间，
	某个其他线程将\co{a}的值设为零。
	当可怜的用户试图单步执行程序时，惊讶地发现
	程序现在在else子句中而不是then子句中。

	这\emph{不是}我所说的易于使用的调试器。
}\QuickQuizEnd

当然，中止和回滚引发了HTM是否可用于硬实时系统的问题。
HTM的性能优势是否超过了中止和回滚的代价，如果是，在什么条件下？
事务可以使用优先级提升吗？
还是高优先级线程的事务应该优先中止低优先级线程的事务？
如果是这样，如何有效地通知硬件优先级？
关于HTM实时使用的文献相当稀少，
也许是因为在非实时环境中使HTM良好工作就已经有足够多的问题了。

因为当前的HTM实现可能确定性地中止给定事务，
软件必须提供后备代码。
这种后备代码必须使用某种其他形式的同步，例如锁。
如果曾经使用过基于锁的后备，那么锁的所有限制，
包括\IX{deadlock}的可能性，都会重新出现。
当然可以希望后备不经常使用，这可能允许使用更简单、
更不容易死锁的锁设计。
但这引发了系统如何从使用基于锁的后备过渡回事务的问题。\footnote{
	应用程序卡在后备模式中的可能性被称为"旅鼠效应"，
	这个术语归功于Dave Dice。}
一种方法是使用test-and-test-and-set策略~\cite{Martinez02a}，
让所有人都等到锁释放，允许系统在那个时点以事务模式从干净的状态开始。
然而，这可能导致相当多的自旋，如果锁持有者已被阻塞或被抢占，
这可能不是明智之举。
另一种方法是允许事务与持有锁的线程并行进行~\cite{Martinez02a}，
但这在维护原子性方面带来了困难，
特别是如果线程持有锁的原因是因为相应的事务无法放入缓存。

最后，处理中止和回滚的可能性似乎给开发者增加了额外的负担，
他们必须正确处理所有可能的错误条件组合。

显然，HTM的用户必须在测试后备代码路径和从后备代码
过渡回事务代码方面投入大量验证工作。
也没有理由认为HTM硬件的验证要求不那么艰巨。

\subsubsection{缺乏前进保证}
\label{sec:future:Lack of Forward-Progress Guarantees}

尽管事务大小、冲突和中止/回滚都可能导致事务中止，
人们可能希望足够小且持续时间短的事务可以被保证最终成功。
这将允许无条件重试事务，就像\IXacrmf{cas}和load-linked/store-conditional
（LL/SC）操作在使用这些指令实现原子操作（atomic operation）的代码中被无条件重试一样。

不幸的是，除了低时钟频率的学术研究原型~\cite{MartinSchoeberl2010realtimeTM}之外，
当前可用的HTM实现拒绝做出任何形式的前进进度保证。
如前所述，HTM因此不能用于在这些系统上避免死锁。
希望未来的HTM实现将提供某种形式的前进进度保证。
在此之前，HTM在实时应用中必须极其谨慎地使用。

截至2021年，这一暗淡图景的唯一例外是IBM大型机，
它提供了\emph{受约束事务}~\cite{ChristianJacobi2012MainframeTM}。
约束条件相当严格，在\cref{sec:future:Forward-Progress Guarantees}中介绍。
看看HTM前进进度保证是否会从大型机迁移到通用CPU系列将是有趣的。

\subsubsection{不可撤销操作}
\label{sec:future:Irrevocable Operations}

中止和回滚的另一个后果是HTM事务无法容纳不可撤销的操作。
当前的HTM实现通常通过要求事务中的所有访问都是对可缓存内存的访问
（从而禁止MMIO访问）并在中断、陷阱和异常时中止事务
（从而禁止系统调用）来执行此限制。

请注意，只要缓冲区填充/刷新操作在事务外发生，
HTM事务可以容纳缓冲I/O。
这之所以有效，是因为向缓冲区添加数据和从中移除数据是可撤销的：
只有实际的缓冲区填充/刷新操作是不可撤销的。
当然，这种缓冲I/O方法的效果是将I/O包含在事务的占用空间中，
增加了事务的大小，从而增加了失败的概率。

\subsubsection{语义差异}
\label{sec:future:Semantic Differences}

虽然HTM在许多情况下可以作为锁的直接替代品使用
（因此得名\IXacrfst{tle}~\cite{DaveDice2008TransactLockElision}），
但语义上存在细微差异。
Blundell~\cite{Blundell2006TMdeadlock}给出了一个特别棘手的例子，
涉及协调的基于锁的临界区，在事务方式执行时会导致死锁或\IX{livelock}，
但一个更简单的例子是空临界区。

在基于锁的程序中，一个空临界区将保证
之前持有该锁的所有进程现在都已释放它。
2.4版Linux内核的网络栈使用了这种习惯用法来协调配置更改。
但如果将这个空临界区转换为事务，结果就是一个空操作。
所有先前临界区已终止的保证将丢失。
换句话说，事务锁省略保留了锁的数据保护语义，
但丢失了锁的基于时间的消息传递语义。

\QuickQuizSeries{%
\QuickQuizB{
	但为什么\emph{任何人}会需要一个空的基于锁的临界区？
}\QuickQuizAnswerB{
	参见\cref{sec:locking:Exclusive Locks}中
	\QuickQuizARef{\QlockingQemptycriticalsection}的答案。

	然而，有人声称，给定一个没有前进进度保证的强原子\IXacr{htm}
	实现，任何基于空临界区的基于内存的锁设计
	在事务锁省略存在的情况下将正确运行。
	虽然我没有看到这个声明的证明，但有一个直接的理由。
	主要思想是，在强原子HTM实现中，
	给定事务的结果在事务成功完成之前是不可见的。
	因此，如果你能看到一个事务已经开始，
	则保证它已经完成，这意味着
	后续的空锁临界区将成功地"等待"它——毕竟，
	不需要等待。

	这种推理不适用于弱原子系统
	（包括许多\IXacr{stm}实现），也不适用于
	使用内存以外的方式进行通信的基于锁的程序。
	这样的方式之一是时间的流逝（例如，在硬实时系统中）
	或优先级的流动（例如，在软实时系统中）。

	依赖优先级提升的锁设计特别值得关注。
}\QuickQuizEndB
%
\QuickQuizM{
	事务锁省略难道不能通过简单地选择不省略空的基于锁的临界区
	来轻松处理锁的基于时间的消息传递语义吗？
}\QuickQuizAnswerM{
	它可以这样做，但这既不必要也不充分。

	在空临界区由条件编译导致的情况下，这是不必要的。
	在这种情况下，锁的唯一目的很可能是保护数据，
	因此完全省略它是正确的做法。
	实际上，保留空的基于锁的临界区会降低性能和可扩展性。

	另一方面，非空的基于锁的临界区可能同时依赖于
	锁的数据保护和基于时间的消息传递语义。
	在这种情况下使用事务锁省略将是不正确的，
	并会导致错误。
}\QuickQuizEndM
%
\QuickQuizE{
	鉴于现代硬件~\cite{PeterOkech2009InherentRandomness}，
	怎么能指望依赖时序的并行软件正常工作呢？
}\QuickQuizAnswerE{
	简短的回答是，在普通的商用硬件上，
	基于任何精细粒度时序的同步设计都是不明智的，
	不能指望在所有条件下正确运行。

	话虽如此，有些专为硬实时使用而设计的系统要确定性得多。
	在你使用这样的系统的（非常不可能的）情况下，
	这里有一个展示基于时间的同步如何工作的简单例子。
	再次强调，\emph{不要}在商用微处理器上尝试这种方法，
	因为它们具有高度不确定性的性能特征。

	这个例子使用多个工作线程和一个控制线程。
	每个工作线程对应一个出站数据流，
	在每个工作单元执行后将当前时间
	（例如，从\co{clock_gettime()}系统调用获取）
	记录在每线程的\co{my_timestamp}变量中。
	此示例的实时性质导致了以下约束集：

	\begin{enumerate}
	\item	如果给定的工作线程未能在超过
		\co{MAX_LOOP_TIME}的时间段内更新其时间戳，
		则为致命错误。
	\item	锁被谨慎地用于访问和更新全局状态。
	\item	锁在给定线程优先级内按严格的FIFO顺序授予。
	\end{enumerate}

	当工作线程完成其数据流时，它们必须与应用程序的其余部分解除关联，
	并在初始化为\co{-1}的每线程\co{my_status}变量中
	放置一个状态值。
	线程不退出；它们被放在线程池中以适应后续的处理需求。
	控制线程根据需要分配（和重新分配）工作线程，
	并维护线程状态的直方图。
	控制线程以不高于工作线程的实时优先级运行。

	工作线程的代码如下：

\begin{VerbatimN}
	int my_status = -1;  /* Thread local. */

	while (continue_working()) {
		enqueue_any_new_work();
		wp = dequeue_work();
		do_work(wp);
		my_timestamp = clock_gettime(...);
	}

	acquire_lock(&departing_thread_lock);

	/*
	 * Disentangle from application, might
	 * acquire other locks, can take much longer
	 * than MAX_LOOP_TIME, especially if many
	 * threads exit concurrently.
	 */
	my_status = get_return_status();
	release_lock(&departing_thread_lock);

	/* thread awaits repurposing. */
\end{VerbatimN}

	控制线程的代码如下：

\begin{fcvlabel}[ln:future:htm:control thread]
\begin{VerbatimN}[commandchars=\\\@\$]
	for (;;) {
		for_each_thread(t) {
			ct = clock_gettime(...);
			d = ct - per_thread(my_timestamp, t);
			if (d >= MAX_LOOP_TIME) {	\lnlbl@if$
				/* thread departing. */	\lnlbl@dep:b$
				acquire_lock(&departing_thread_lock); \lnlbl@acq$
				release_lock(&departing_thread_lock); \lnlbl@rel$
				i = per_thread(my_status, t);
				status_hist[i]++; /* Bug if TLE! */ \lnlbl@dep:e$
			}
		}
		/* Repurpose threads as needed. */
	}
\end{VerbatimN}
\end{fcvlabel}

	\begin{fcvref}[ln:future:htm:control thread]
	\Clnref{if}使用时间的流逝来推断线程已退出，
	如果是这样则执行\clnref{dep:b,dep:e}。
	\clnref{acq,rel}上的空锁临界区保证任何正在退出过程中的线程
	已经完成（记住锁是按FIFO顺序授予的！）。
	\end{fcvref}

	再次强调，不要在商用微处理器上尝试这种方法。
	毕竟，即使在专门为硬实时使用而设计的系统上，
	正确实现这一点也够困难的了！
}\QuickQuizEndE
}

锁和事务之间的一个重要语义差异是在基于锁的实时程序中
用于避免优先级反转的优先级提升。
优先级反转可能发生的一种方式是当持有锁的低优先级线程
被中等优先级的CPU密集型线程抢占时。
如果每个CPU至少有一个这样的中等优先级线程，
低优先级线程将永远没有机会运行。
如果高优先级线程现在试图获取该锁，它将被阻塞。
它无法获取锁直到低优先级线程释放它，
低优先级线程无法释放锁直到它有机会运行，
而它无法有机会运行直到某个中等优先级线程放弃其CPU。
因此，中等优先级线程实际上阻塞了高优先级进程，
这就是"优先级反转"这个名称的由来。

\begin{listing}
\begin{fcvlabel}[ln:future:Exploiting Priority Boosting]
\begin{VerbatimL}[commandchars=\\\@\$]
void boostee(void)		\lnlbl@low:b$
{
	int i = 0;

	acquire_lock(&boost_lock[i]);	\lnlbl@1stacq$
	for (;;) {
		acquire_lock(&boost_lock[!i]);
		release_lock(&boost_lock[i]);
		i = i ^ 1;
		do_something();
	}
}				\lnlbl@low:e$

void booster(void)		\lnlbl@high:b$
{
	int i = 0;

	for (;;) {
		usleep(500); /* sleep 0.5 ms. */
		acquire_lock(&boost_lock[i]);	\lnlbl@acq$
		release_lock(&boost_lock[i]);	\lnlbl@rel$
		i = i ^ 1;
	}
}                               \lnlbl@high:e$
\end{VerbatimL}
\end{fcvlabel}
\caption{利用优先级提升}
\label{lst:future:Exploiting Priority Boosting}
\end{listing}

避免优先级反转的一种方法是\emph{优先级继承}，
其中被锁阻塞的高优先级线程临时将其优先级捐赠给锁的持有者，
这也被称为\emph{优先级提升}。
然而，优先级提升可以用于避免优先级反转以外的目的，
如\cref{lst:future:Exploiting Priority Boosting}所示。
\begin{fcvref}[ln:future:Exploiting Priority Boosting]
该清单的\clnrefrange{low:b}{low:e}展示了一个必须大约每毫秒运行一次的低优先级进程，
而同一清单的\clnrefrange{high:b}{high:e}展示了一个使用优先级提升
来确保\co{boostee()}按需定期运行的高优先级进程。

\co{boostee()}函数通过始终持有两个\co{boost_lock[]}锁中的一个
来实现这一点，以便\co{booster()}的\clnrefrange{acq}{rel}
可以根据需要提升优先级。
\end{fcvref}

\QuickQuiz{
	但\cref{lst:future:Exploiting Priority Boosting}中的
	\co{boostee()}函数交替以相反的顺序获取锁！
	这不会导致死锁吗？
}\QuickQuizAnswer{
	不会导致死锁。
	要产生死锁，两个不同的线程必须各自以相反的顺序
	获取两个锁，这在此示例中不会发生。
	然而，像lockdep~\cite{JonathanCorbet2006lockdep}
	这样的死锁检测器会将此标记为误报。
}\QuickQuizEnd

\begin{fcvref}[ln:future:Exploiting Priority Boosting]
这种安排要求\co{boostee()}在系统变忙之前在\clnref{1stacq}处获取其第一个锁，
但即使在现代硬件上也很容易安排。

不幸的是，这种安排在事务锁省略存在的情况下可能会失效。
\co{boostee()}函数重叠的临界区变成一个无限事务，
迟早会中止，例如在运行\co{boostee()}函数的线程
第一次被抢占时。
此时\co{boostee()}将回退到使用锁，但鉴于其低优先级
以及安静的初始化期间已经结束
（这毕竟是\co{boostee()}被抢占的原因），
这个线程可能永远不会再有机会运行。

而如果\co{boostee()}线程没有持有锁，那么
\cref{lst:future:Exploiting Priority Boosting}中
\co{booster()}线程在\clnref{acq,rel}上的空临界区
将变成一个没有效果的空事务，
从而\co{boostee()}永远不会运行。
此例说明了transactional memory的回滚重试语义的一些微妙后果。
\end{fcvref}

鉴于经验可能会揭示更多微妙的语义差异，
基于HTM的锁省略在大型程序中的应用应当谨慎进行。
话虽如此，在适用的地方，基于HTM的锁省略可以消除
与锁变量相关的缓存未命中，这在截至2015年初的
大型实际软件系统中已经带来了数十个百分点的性能提升。
因此我们可以期望在提供可靠支持的硬件上看到这种技术的大量使用。

\QuickQuiz{
	所以一群人着手取代锁，结果大多只是在优化锁？
}\QuickQuizAnswer{
	至少他们完成了一些有用的事情！
	也许随着时间的推移，\IXacr{htm}会继续取得进步~\cite{Siakavaras2017CombiningHA,DimitriosSiakavaras2020RCU-HTM-B+Trees,ChristinaGiannoula2018HTM-RCU-graphcoloring,SeongJaePark2020HTMRCUlock}。
}\QuickQuizEnd

\subsubsection{总结}
\label{sec:future:HTM Weaknesses WRT Locking: Summary}

% future/HTMtable.tex
% SPDX-License-Identifier: CC-BY-SA-3.0

\begin{table*}
\centering
% \scriptsize
\small
\input{future/HTMtableColor}
\setlength{\tabcolsep}{4pt}\OneColumnHSpace{-.9in}
\ebresizewidth{
\begin{tabularx}{6.5in}{p{0.95in}cXcX}
\toprule
  &
    & \multicolumn{1}{c}{锁} &
      & \multicolumn{1}{c}{Hardware Transactional Memory} \\
\midrule
  基本思想 &
    & 一次只允许一个线程访问给定的一组对象。 &
      & 使对给定对象集的操作以原子方式执行。 \\
\midrule
  适用范围 &
    & \Pl 处理所有操作。 &
      & \Pl 处理可撤销操作。 \\
\addlinespace[4pt]
  &
    & &
      & \Mn 不可撤销操作迫使回退（通常回退到锁）。 \\
\midrule
  可组合性 &
    & \Dw 受死锁限制。 &
      & \Dw 受不可撤销操作、事务大小和死锁限制
        （假设基于锁的回退代码）。 \\
\midrule
  可扩展性与性能 &
    & \Mn 数据必须可分区以避免锁竞争。 &
      & \Mn 数据必须可分区以避免冲突。 \\
\cmidrule{3-5}
  &
    & \Dw 分区通常必须在设计时固定。 &
      & \Pl 自动动态调整分区，精确到缓存行边界。 \\
\addlinespace[4pt]
  &
    & &
      & \Mn 回退需要分区（对于罕见的回退不太重要）。 \\
\cmidrule{3-5}
  &
    & \Dw 锁原语通常导致昂贵的缓存未命中
      和内存屏障指令。 &
      & \Mn 事务开始/结束指令通常不会导致缓存未命中，
        但确实有内存排序和开销影响。 \\
\cmidrule{3-5}
  &
    & \Pl 竞争效应集中在获取和释放上，因此
      临界区以全速运行。 &
      & \Mn 竞争会中止冲突的事务，即使它们已经
        运行了很长时间。 \\
\cmidrule{3-5}
  &
    & \Pl 私有化操作简单、直观、高效且可扩展。 &
      & \Mn 私有化数据增加事务大小。 \\
\midrule
  硬件支持 &
    & \Pl 通用硬件即可。 &
      & \Mn 需要新硬件（正开始变得可用）。 \\
\cmidrule{3-5}
  &
    & \Pl 性能不受缓存几何细节影响。 &
      & \Mn 性能严重依赖缓存几何结构。 \\
\midrule
  软件支持 &
    & \Pl API已存在，有大量代码和经验，调试器正常工作。 &
      & \Mn API正在出现，DBMS之外经验很少，事务中的
        断点可能有问题。 \\
\midrule
  与其他机制的交互 &
    & \Pl 有长期成功交互的经验。 &
      & \Dw 刚刚开始研究交互。 \\
\midrule
  实际应用 &
    & \Pl 有。 &
      & \Pl 有。 \\
\midrule
  广泛适用性 &
    & \Pl 有。 &
      & \Mn 尚无定论。 \\
\bottomrule
\end{tabularx}
}
\IfTwoColumn{
\caption{锁与HTM的比较（\colorbox{plus}{优势}、
  \colorbox{minus}{劣势}、\colorbox{down}{严重劣势}）}
}{
\caption{锁与HTM的比较（\colorbox{plus}{优势}、
  \colorbox{minus}{劣势}、
  \colorbox{down}{严重} \colorbox{down}{劣势}）}
}
\label{tab:future:Comparison of Locking and HTM}
\end{table*}


虽然HTM似乎可能会有引人注目的用例，
但当前的实现存在严重的事务大小限制、冲突处理复杂性、
中止和回滚问题以及需要仔细处理的语义差异。
HTM相对于锁的当前状况总结在
\cref{tab:future:Comparison of Locking and HTM}中。
可以看出，虽然HTM的当前状态缓解了锁的一些严重不足，\footnote{
	公平地说，重要的是要强调锁的不足确实有广为人知且广泛使用的工程解决方案，
	包括死锁检测器~\cite{JonathanCorbet2006lockdep}、
	大量已适应锁的数据结构，以及如
	\cref{sec:future:HTM Weaknesses WRT Locking When Augmented}中讨论的
	长期增强历史。
	此外，如果锁真的像快速浏览许多学术论文可能合理地使人相信的那样糟糕，
	那些大型的基于锁的并行程序（包括自由开源软件和专有软件）
	到底是从哪里来的呢？}
但它通过引入大量自身的不足来实现这一点。
这些不足被TM社区的领导者所承认~\cite{AlexanderMatveev2012PessimisticTM}。\footnote{
	此外，在2011年初，我受邀对transactional memory背后的
	一些假设进行批评~\cite{PaulEMcKenney2011Verico}。
	听众出奇地不敌视，不过也许他们对我手下留情是因为
	我在做报告时严重倒时差。}

此外，这还不是完整的故事。
锁通常不是单独使用的，而是通常由其他同步机制增强，
包括引用计数、原子操作、非阻塞数据结构、
\IXpl{hazard pointer}~\cite{MagedMichael04a,HerlihyLM02}
和RCU~\cite{McKenney98,McKenney01a,ThomasEHart2007a,PaulEMcKenney2012ELCbattery}。
下一节将探讨这种增强如何改变等式。

\subsection{增强后的 HTM 相对于锁的劣势}
\label{sec:future:HTM Weaknesses WRT Locking When Augmented}

% future/HTMtableRCU.tex
% SPDX-License-Identifier: CC-BY-SA-3.0

\begin{table*}
\centering
\small
\input{future/HTMtableColor}
\setstretch{0.95}
\setlength{\tabcolsep}{4pt}\OneColumnHSpace{-.9in}
\ebresizewidth{
\resizebox{6.5in}{!}{
\begin{tabularx}{6.8in}{p{0.95in}cXcX}
\toprule
  &
    & \multicolumn{1}{c}{用户空间RCU或Hazard Pointer增强的锁} &
      & \multicolumn{1}{c}{Hardware Transactional Memory} \\
\midrule
  基本思想 &
    & 一次只允许一个线程访问给定的一组对象。 &
      & 使对给定对象集的操作以原子方式执行。 \\
\midrule
  适用范围 &
    & \Pl 处理所有操作。 &
      & \Pl 处理可撤销操作。 \\
\addlinespace[4pt]
  &
    & &
      & \Mn 不可撤销操作迫使回退（通常回退到锁）。 \\
\midrule
  可组合性 &
    & \Pl 读取端仅受宽限期等待操作限制。 &
      & \Dw 受不可撤销操作、事务大小和死锁限制
        （假设基于锁的回退代码）。\\
\addlinespace[4pt]
  &
    & \Mn 更新端受死锁限制。读取端减少死锁。 &
      & \\
\midrule
  可扩展性与性能 &
    & \Mn 数据必须可分区以避免更新端之间的锁竞争。 &
      & \Mn 数据必须可分区以避免冲突。 \\
\addlinespace[4pt]
  &
    & \Pl 读取端不需要分区。 &
      & \\
\cmidrule{3-5}
  &
    & \Dw 更新端的分区通常必须在设计时固定。 &
      & \Pl 自动动态调整分区，精确到缓存行边界。 \\
\cmidrule{3-5}
  &
    & \Pl 读取端不需要分区。 &
      & \Mn 回退需要分区（对于罕见的回退不太重要）。 \\
\cmidrule{3-5}
  &
    & \Dw 更新端锁原语通常导致昂贵的缓存未命中
      和内存屏障指令。 &
      & \Mn 事务开始/结束指令通常不会导致缓存未命中，
        但确实有内存排序和开销影响。 \\
\cmidrule{3-5}
  &
    & \Pl 更新端竞争效应集中在获取和释放上，
      因此临界区以全速运行。 &
      & \Mn 竞争会中止冲突的事务，即使它们已经
        运行了很长时间。 \\
\addlinespace[4pt]
  &
    & \Pl 读取端不与更新端或彼此竞争。 &
      & \\
\cmidrule{3-5}
  &
    & \Pl 读取端原语通常是有界无等待的，开销低。
      （Hazard pointer为无锁，开销低。） &
      & \Mn 只读事务受冲突和回滚影响。
        除回退代码提供的之外没有前向进展保证。 \\
\cmidrule{3-5}
  &
    & \Pl 当数据仅对更新端可见时，私有化操作
      简单、直观、高效且可扩展。 &
      & \Mn 私有化数据增加事务大小。 \\
\addlinespace[4pt]
  &
    & \Mn 对于读取端可见的数据，私有化操作开销较大
      （但仍然直观且可扩展）。 &
      & \\
\midrule
  硬件支持 &
    & \Pl 通用硬件即可。 &
      & \Mn 需要新硬件（正开始变得可用）。 \\
\cmidrule{3-5}
  &
    & \Pl 性能不受缓存几何细节影响。 &
      & \Mn 性能严重依赖缓存几何结构。 \\
\midrule
  软件支持 &
    & \Pl API已存在，有大量代码和经验，调试器正常工作。 &
      & \Mn API正在出现，DBMS之外经验很少，事务中的
        断点可能有问题。 \\
\midrule
  与其他机制的交互 &
    & \Pl 有长期成功交互的经验。 &
      & \Dw 刚刚开始研究交互。 \\
\midrule
  实际应用 &
    & \Pl 有。 &
      & \Pl 有。 \\
\midrule
  广泛适用性 &
    & \Pl 有。 &
      & \Mn 尚无定论。 \\
\bottomrule
\end{tabularx}
}}
\caption{锁（由RCU或Hazard Pointer增强）与HTM的比较
  （\colorbox{plus}{优势}、\colorbox{minus}{劣势}、
  \colorbox{down}{严重劣势}）}
\label{tab:future:Comparison of Locking (Augmented by RCU or Hazard Pointers) and HTM}
\end{table*}


实践者长期以来使用引用计数、原子操作、非阻塞数据结构、
hazard pointer和RCU来避免锁的一些不足。
例如，在许多情况下，可以通过使用引用计数、hazard pointer或RCU
来保护数据结构来避免死锁，
特别是对于只读临界区~\cite{MagedMichael04a,HerlihyLM02,MathieuDesnoyers2012URCU,DinakarGuniguntala2008IBMSysJ,ThomasEHart2007a}。
这些方法还减少了划分数据结构的需要，
正如在\cref{chp:Data Structures}中所见。
RCU进一步提供了无竞争的有界无等待读端原语~\cite{McKenney98,MathieuDesnoyers2012URCU}，
而hazard pointer提供了无锁的读端原语~\cite{Michael02a,HerlihyLM02,MagedMichael04a}。
将这些考虑因素添加到\cref{tab:future:Comparison of Locking and HTM}中，
得到了\cref{tab:future:Comparison of Locking (Augmented by RCU or Hazard Pointers) and HTM}
中所示的增强锁与HTM之间的更新比较。
两个表之间差异的总结如下：

\begin{enumerate}
\item	使用非阻塞读端机制缓解了死锁问题。
\item	hazard pointer和RCU等读端机制可以在不可划分的数据上高效操作。
\item	hazard pointer和RCU不会彼此竞争，也不会与更新者竞争，
	允许读多写少工作负载获得出色的性能和可扩展性。
\item	hazard pointer和RCU提供前进进度保证
	（分别为无锁和有界无等待）。
\item	hazard pointer和RCU的私有化操作是直接的。
\end{enumerate}

\IfEbookSize{}{
\input{future/HTMtableFull}

对于视力好的人，
\cref{tab:future:Comparison of Locking (Plain and Augmented) and HTM}
合并了\cref{tab:future:Comparison of Locking and HTM,%
tab:future:Comparison of Locking (Augmented by RCU or Hazard Pointers) and HTM}。
}

\QuickQuiz{
	\Cref{tab:future:Comparison of Locking and HTM,tab:future:Comparison of Locking (Augmented by RCU or Hazard Pointers) and HTM}
	声称硬件才刚刚开始可用。
	但HTM硬件支持不是已经广泛可用近十年了吗？
}\QuickQuizAnswer{
	是也不是。
	看来即使是在真实硬件中实现TM的HTM子集也比看起来要复杂一些~\cite{ChristianJacobi2012MainframeTM,ScottWasson2014HaswellDisableTSX,Intel2020HaswellTSXordering,Intel2021perfTSXordering,MichaelLarabel2021DisableTSX}。
	因此，令人遗憾的事实是，截至2021年"刚刚开始可用"
	这一说法过于准确了。
	实际上，供应商已经开始弃用其HTM实现~\cite[Book III Appendix A]{PowerISA3.1-2020}。
}\QuickQuizEnd

当然，也可以增强HTM，如下一节中讨论的那样。

\subsection{HTM 最适合用在哪里？}
\label{sec:future:Where Does HTM Best Fit In?}

虽然HTM的适用范围可能还需要一段时间才能像
\cpageref{fig:defer:RCU Areas of Applicability}上
\cref{fig:defer:RCU Areas of Applicability}中所示的RCU那样清晰地划定，
但这不是不开始朝这个方向前进的理由。

HTM似乎最适合涉及对在大型多处理器上运行的相对大型内存数据结构的
不同部分进行相对小的更改的更新密集型工作负载，
因为这满足了当前HTM实现的大小限制，同时最小化了冲突
及伴随的中止和回滚的概率。
这种场景也是给定当前同步原语相对难以处理的场景。

将锁与HTM结合使用似乎可以克服HTM在不可撤销操作方面的困难，
而使用RCU或hazard pointer可能会缓解HTM对遍历数据结构大部分内容的
只读操作的事务大小限制~\cite{SeongJaePark2020HTMRCUlock}。
当前的HTM实现会无条件中止与RCU或hazard pointer读者冲突的更新事务，
但也许未来的HTM实现将与这些同步机制更顺畅地互操作。
与此同时，更新与大型RCU或hazard pointer读端临界区冲突的概率
应该远小于与等效只读事务冲突的概率。\footnote{
	相当讽刺的是，严格的事务机制出现在共享内存系统中的时间
	恰好是NoSQL数据库放松传统数据库应用对严格事务依赖的时候。
	尽管如此，HTM确实实现了TM的易用性承诺，
	尽管是用于对Linux内核地址空间随机化防御机制的
	黑帽攻击~\cite{YeongjinJang2016TSXbreakKASLR,Jang:2016:BKA:2976749.2978321}。}
尽管如此，持续不断的RCU或hazard pointer读者完全有可能
因为相应的持续冲突流而使更新者饥饿。
可以通过（以显著的硬件成本和复杂性为代价）
给事务外读取提供被加载内存位置的事务前副本来消除这一漏洞。

HTM事务必须有后备方案的事实可能在某些情况下
会将数据结构的静态可划分性重新强加给HTM。
如果未来的HTM实现提供前进进度保证，
这一限制可能会得到缓解，这可能在某些情况下消除对后备代码的需要，
进而可能允许HTM在冲突概率较高的情况下高效使用。

简而言之，虽然HTM可能有重要的用途和应用，
但它是并行程序员工具箱中的另一个工具，而不是工具箱整体的替代品。

\subsection{潜在的变革因素}
\label{sec:future:Potential Game Changers}

可能大大增加HTM需求的变革因素包括以下几点：

\begin{enumerate}
\item	前进进度保证。
\item	事务大小增加。
\item	改进的调试支持。
\item	弱原子性。
\end{enumerate}

以下各节对此进行了展开讨论。

\subsubsection{前进保证}
\label{sec:future:Forward-Progress Guarantees}

正如\cref{sec:future:Lack of Forward-Progress Guarantees}中所讨论的，
当前的HTM实现缺乏前进保证，这要求必须有回退软件来处理HTM失败。
当然，提出保证要求很容易，但提供保证并不总是那么容易。
对于HTM而言，实现保证的障碍包括缓存大小和关联度、TLB大小和关联度、
事务持续时间和中断频率，以及调度器的实现。

缓存大小和关联度已在\cref{sec:future:Transaction-Size Limitations}中讨论过，
同时还介绍了一些旨在解决当前限制的研究工作。
然而，HTM的前进保证将伴随着大小限制，尽管这些限制可能有朝一日会变得很大。
那么，为什么当前的HTM实现不对小事务提供前进保证呢？例如，将事务限制在
缓存的关联度范围内？
一个可能的原因是需要处理硬件故障。
例如，一个出故障的缓存SRAM单元可能通过停用该故障单元来处理，
从而降低缓存的关联度，因此也降低了能够保证前进的事务的最大尺寸。
鉴于这只是减少了保证的事务大小，看来可能还有其他原因在起作用。
也许在生产级硬件上提供前进保证比人们想象的更加困难，
考虑到在软件中做出前进保证的难度，这是一个完全合理的解释。
将问题从软件转移到硬件并不一定使其更容易解决~\cite{ChristianJacobi2012MainframeTM}。

对于使用物理标记和物理索引的缓存，事务能够放入缓存是不够的。
它的地址转换还必须能放入TLB\@。
因此，任何前进保证还必须考虑TLB的大小和关联度。

鉴于在当前的HTM实现中，中断、陷阱和异常都会中止事务，
因此给定事务的执行持续时间必须短于预期的中断间隔。
无论给定事务接触的数据多么少，如果它运行时间过长，就会被中止。
因此，任何前进保证不仅必须以事务大小为条件，还必须以事务持续时间为条件。

前进保证在很大程度上取决于确定多个冲突事务中哪个应该被中止的能力。
很容易想象一个无尽的事务系列，每个事务中止一个更早的事务，
而自己又被一个更晚的事务中止，以至于没有任何事务能真正提交。
冲突处理的复杂性从已经提出的大量HTM冲突解决策略中可见一斑~\cite{EgeAkpinar2011HTM2TLE,YujieLiu2011ToxicTransactions}。
事务外的访问引入了额外的复杂性，
正如Blundell~\cite{Blundell2006TMdeadlock}所指出的那样。
很容易将所有这些问题归咎于事务外的访问，
但这种思路的荒谬之处可以通过将每个事务外的访问放入其自身的
单次访问事务中来轻松证明。
问题在于访问模式本身，而不是它们是否恰好被封装在事务中。

最后，事务的任何前进保证还取决于调度器，
它必须让执行事务的thread运行足够长的时间来成功提交。

因此，HTM厂商提供前进保证存在重大障碍。
然而，如果其中任何一家做到了，其影响将是巨大的。
这意味着HTM事务将不再需要软件回退，
也就意味着HTM终于可以兑现TM消除死锁的承诺。

然而，在2012年末，IBM大型机宣布了一种HTM实现，
除了通常的尽力而为HTM实现外，还包括\emph{受约束事务}~\cite{ChristianJacobi2012MainframeTM}。
受约束事务以\co{tbeginc}指令开始，
而不是用于尽力而为事务的\co{tbegin}指令。
受约束事务保证总是会（最终）完成，
所以如果事务中止，硬件不是跳转到回退路径
（如同尽力而为事务那样），而是在\co{tbeginc}指令处重新启动事务。

大型机架构师需要采取极端措施来兑现这个前进保证。
如果给定的受约束事务反复失败，CPU可能会禁用分支预测、
强制顺序执行，甚至禁用流水线。
如果反复失败是由于高竞争造成的，CPU可能会禁用推测性取指、
引入随机延迟，甚至序列化冲突CPU的执行。
"有趣的"前进场景涉及少则两个CPU，多则一百个CPU。
也许这些极端措施提供了一些线索，解释了为什么其他CPU
至今没有提供受约束事务。

顾名思义，受约束事务实际上受到严格的约束：

\begin{enumerate}
\item	最大数据占用量是四个内存块，
	其中每个块不能大于32字节。
\item	最大代码占用量是256字节。
\item	如果给定的4K页面包含受约束事务的代码，
	则该页面不能包含该事务的数据。
\item	可执行的最大汇编指令数是32条。
\item	禁止向后跳转。
\end{enumerate}

尽管如此，这些约束支持许多重要的数据结构，
包括链表、栈、队列和数组。
因此，受约束HTM似乎有可能成为并行程序员工具箱中的重要工具。

请注意，这些前进保证不必是绝对的。
例如，假设一种HTM的使用以全局锁作为回退。
假设回退机制已经被仔细设计以避免
\cref{sec:future:Aborts and Rollbacks}中讨论的"旅鼠效应"，
那么如果HTM回滚足够不频繁，全局锁就不会成为瓶颈。
话虽如此，系统越大、临界区越长、
从"旅鼠效应"中恢复所需的时间越长，
"足够不频繁"就需要越罕见。

\subsubsection{事务大小增加}
\label{sec:future:Transaction-Size Increases}

前进保证很重要，但正如我们所看到的，它们将是基于事务大小和持续时间的
有条件保证。
已经取得了一些进展，例如，一些商用HTM实现使用近似技术来支持
极其庞大的HTM读集~\cite{RaviRajwar2012TSX}。
再例如，\Power{8} HTM支持挂起事务，
这可以避免将不相关的访问添加到被挂起事务的读集和
写集中~\cite{Le:2015:TMS:3266491.3266500}。
这一能力已被用于产生高性能的
读写锁~\cite{PascalFelber2016rwlockElision}。

需要注意的是，即使是小尺寸的保证也将非常有用。
例如，两条cache line的保证就足以实现栈、队列或双端队列。
然而，更大的数据结构需要更大的保证，例如，
按顺序遍历一棵树需要与树中节点数量相等的保证。
因此，即使是适度增加保证的大小也会增加HTM的实用性，
从而增加CPU提供这种保证或提供足够好的替代方案的需求。

\subsubsection{改进的调试支持}
\label{sec:future:Improved Debugging Support}

事务大小的另一个抑制因素是调试事务的需求。
当前机制的问题在于，单步异常会中止外层的事务。
对这个问题有许多变通方法，包括模拟处理器（慢！）、
用STM替代HTM（慢而且语义略有不同！）、
使用重复重试来模拟前进的回放技术（奇怪的失败模式！），以及
完全支持调试HTM事务（复杂！）。

如果某个HTM厂商能够生产出允许在事务内直接使用经典调试技术的
HTM系统，包括断点、单步执行和打印语句，
这将使HTM更具吸引力。
一些事务内存（transactional memory）研究者在2013年开始认识到这个问题，
至少有一个提案涉及硬件辅助调试设施~\cite{JustinGottschlich2013TMdebug}。
当然，这个提案依赖于现成的硬件获得这样的
设施~\cite{TimothyHayes2020ARM-HTM,Intel2020TSXdevguide}。
更糟糕的是，一些前沿的调试设施与HTM不兼容~\cite{RobertOCallahan2020DebuggingHTM}。

\subsubsection{弱原子性}
\label{sec:future:Weak Atomicity}

鉴于HTM在可预见的未来可能面临某种大小限制，
HTM有必要与其他机制顺畅地互操作。
如果事务外的读取不会无条件地中止具有冲突写入的事务——
而是简单地为读取提供事务前的值——
那么HTM与风险指针（hazard pointer）和RCU等以读为主的机制之间的互操作性将会得到改善。
这样，风险指针和RCU可以被用来允许HTM处理更大的数据结构并减少冲突概率。

然而，这并不一定简单。
最直接的实现方式需要在每条cache line和总线上增加一个额外的状态，
这是一笔不小的额外开销。
与这笔开销相伴的好处是允许大占用量的读取者，
而不会因持续冲突而饿死更新者。
一种替代方法由Siakavaras等人~\cite{Siakavaras2017CombiningHA}
对二叉搜索树取得了极好的效果，
即使用RCU进行只读遍历，而仅使用HTM进行实际的更新操作。
这种组合的性能比其他事务内存技术高出多达220\pct，
与Howard和Walpole~\cite{PhilHoward2011RCUTMRBTree}
将RCU与STM结合时观察到的加速比相似\@。
在这两种情况下，弱原子性都是在软件中而非硬件中实现的。
尽管如此，看看在硬件和软件中同时实现弱原子性能获得
哪些额外的加速，仍然是一件有趣的事情。

\subsection{结论}
\label{sec:future:Conclusions}

虽然当前的HTM实现在某些情况下确实带来了真正的性能优势，
但它们也有显著的缺陷。
最显著的缺陷似乎是
有限的事务大小、
冲突处理的需求、中止和回滚的需求、
缺乏前进保证、
无法处理不可撤销操作，
以及与加锁之间微妙的语义差异。
关于HTM实现可靠性，也有挥之不去的
担忧~\cite{ChristianJacobi2012MainframeTM,ScottWasson2014HaswellDisableTSX,Intel2020HaswellTSXordering,Intel2021perfTSXordering,MichaelLarabel2021DisableTSX,PowerISA3.1-2020}。

其中一些缺陷可能在未来的实现中得到缓解，
但似乎仍然存在强烈的需求，使HTM能够与许多其他类型的同步机制良好配合，
正如之前所指出的~\cite{McKenney2007PLOSTM,PaulEMcKenney2010OSRGrassGreener}。
虽然已经有一些将HTM与RCU结合使用的
工作~\cite{Siakavaras2017CombiningHA,DimitriosSiakavaras2020RCU-HTM-B+Trees,ChristinaGiannoula2018HTM-RCU-graphcoloring,SeongJaePark2020HTMRCUlock}，
但几乎没有证据表明HTM在与RCU以及其他延迟回收机制的配合方面取得了进展。

简而言之，当前的HTM实现似乎是并行程序员工具箱中受欢迎且有用的补充，
要利用好它们还需要大量有趣且具有挑战性的工作。
然而，它们不能被视为一根魔杖，
用来消除所有并行编程问题。

\QuickQuiz{
	但随着持续的工作，HTM最终实现完整的TM愿景不是必然的吗？
}\QuickQuizAnswer{
	也许吧？

	请回顾\cref{chp:Hardware and its Habits}，尤其是
	对\cref{tab:cpu:CPU 0 View of Synchronization Mechanisms on 8-Socket System With Intel Xeon Platinum 8176 CPUs at 2.10GHz}
	的讨论，它概述了细粒度全局一致的代价。
	还请回顾对\cref{fig:defer:Pre-BSD Routing Table Protected by RCU QSBR}
	的讨论，它展示了避免全局一致需求所带来的可扩展性优势。

	TM主要用于细粒度同步，
	对于每组并发事务，它要求以下全局一致：
	\begin{enumerate*}[(1)]
	\item 事务是有序的，
	或者
	\item 事务之间不冲突。
	\end{enumerate*}
	当然，并发组中可能只有部分事务冲突，
	在这种情况下只有那些冲突的事务需要排序。
	因此，\cref{chp:Hardware and its Habits}中涵盖的内容
	以及对\cref{fig:defer:Pre-BSD Routing Table Protected by RCU QSBR}
	的讨论都直接适用于TM。

	当然，非常聪明的人们确实在继续研究事务内存，
	可以期待他们继续取得进展。
	然而，TM对全局一致的要求与\cref{chp:Hardware and its Habits}中
	讨论的物理定律相冲突，
	这种冲突可能会继续对一般情况下的TM性能和可扩展性施加硬性限制。
}\QuickQuizEnd

% future/formalregress.tex
% mainfile: ../perfbook.tex
% SPDX-License-Identifier: CC-BY-SA-3.0

\section{形式化回归测试？}
\label{sec:future:Formal Regression Testing?}
%
\epigraph{没有实验的理论：
	  我们走得太远了吗？}
	 {Michael Mitzenmacher}

形式化验证长期以来在许多生产环境中证明了其
有用性~\cite{JamesRLarus2004RightingSoftware,AlBessey2010BillionLoCLater,ByronCook2018FormalAmazon,CaitlinSadowski2018staticAnalysisGoogle,DinoDistefano2019FBstaticAnalysis}。
然而，硬核形式化验证是否能被纳入复杂并发代码库（如Linux内核）
的持续集成中所使用的自动化回归测试套件，仍是一个悬而未决的问题。
虽然已经有了Linux内核SRCU的概念验证~\cite{LanceRoy2017CBMC-SRCU}，
但该测试针对的是最简单的RCU实现之一的一小部分，
而且已经证明很难使其跟上不断变化的Linux内核。
因此，值得探讨将形式化验证作为Linux内核回归测试的一等成员
所需的条件。

以下列表是一个好的
起点~\cite[slide 34]{PaulEMcKenney2015DagstuhlVerification}：

\begin{enumerate}
\item	任何所需的转换必须是自动化的。
\item	环境（包括内存排序）必须被正确处理。
\item	内存和CPU开销必须在可接受的适度范围内。
\item	必须提供指向bug位置的具体信息。
\item	除源代码和输入之外的信息范围必须适度。
\item	所定位的bug必须与代码用户相关。
\end{enumerate}

这个列表基于但比Richard Bornat的格言更为温和：
``形式化验证研究人员应该验证开发人员编写的代码，
使用他们编写时所用的语言，在代码运行的环境中运行，
在他们编写时就进行验证。''
以下各节讨论上述每项要求，
然后有一节展示几个工具在这些要求方面表现的记分卡。

\QuickQuiz{
	这个列表太乌托邦了！
	为什么不坚持形式化验证的现有技术水平？
}\QuickQuizAnswer{
	你有权对什么是乌托邦、什么不是乌托邦发表意见，
	但比起任何可能反对这些条目的人，
	我会更关注那些真正在这些条目上取得进展的人。
	这可能与我长期以来的经验有关，
	即有人试图说服我放弃其所偏爱的工具无法处理的特定事项。

	与此同时，请随意阅读那些真正取得进展的人所写的论文，
	例如这篇~\cite{DinoDistefano2019FBstaticAnalysis}。
}\QuickQuizEnd

\subsection{自动转换}
\label{sec:future:Automatic Translation}

虽然Promela和\co{spin}是无价的设计辅助工具，
但如果你需要对C语言程序进行形式化回归测试，
你必须在每次想要重新验证代码时手动将其翻译为Promela。
如果你的代码恰好在Linux内核中（每60--90天发布一次），
你将需要每年手动翻译四到六次。
随着时间推移，人为错误将悄然出现，
这意味着验证将不再与源代码匹配，使验证变得毫无用处。
重复验证显然要求形式化验证工具直接输入你的代码，
或者存在无bug的自动转换将你的代码转换为验证所需的形式。

PPCMEM和\co{herd}理论上可以直接输入汇编语言和C++代码，
但这些工具只能处理非常小的litmus测试，
这通常意味着你必须手动提取你的机制的核心部分。
与Promela和\co{spin}一样，PPCMEM和\co{herd}都极为有用，
但它们不太适合回归测试套件。

相比之下，\IXaltacr{\co{cbmc}}{cbmc}和\IX{Nidhugg}可以输入
合理大小（虽然仍然相当有限）的C程序，
如果它们的能力继续增长，完全可以成为回归测试套件的出色补充。
Coverity静态分析工具也输入C程序，而且可以处理非常大的程序，
包括Linux内核。
当然，Coverity的静态分析与\co{cbmc}和Nidhugg相比要简单得多。
另一方面，Coverity拥有一个包罗万象的``C程序''定义，
这带来了特殊的挑战~\cite{AlBessey2010BillionLoCLater}。
Amazon Web Services使用各种形式化验证工具，
包括\co{cbmc}，并将其中一些工具应用于
回归测试~\cite{ByronCook2018FormalAmazon}。
Google在大型Java代码库上直接使用许多相对简单的静态分析工具，
这些代码库可以说比C代码库更为同质~\cite{CaitlinSadowski2018staticAnalysisGoogle}。
Facebook对其代码库使用更激进的形式化验证形式，
包括并发分析~\cite{DinoDistefano2019FBstaticAnalysis,PeterWOHearn2019incorrectnessLogic}，
尽管尚未应用于Linux内核。
最后，Microsoft长期以来一直对其代码库使用
静态分析~\cite{JamesRLarus2004RightingSoftware}。

鉴于这份清单，显然可以创建直接消费生产级别源代码的
复杂形式化验证工具。

然而，以C代码作为输入的一个不足之处在于它假设编译器是正确的。
另一种方法是将C编译器生成的二进制文件作为输入，
从而考虑到任何相关的编译器bug。
这种方法已在许多验证工作中使用，
最著名的可能是seL4项目~\cite{ThomasSewell2013L4binaryVerification}。

\QuickQuiz{
	鉴于seL4项目中使用的各种验证器的开创性，
	为什么本章不更深入地介绍它们？
}\QuickQuizAnswer{
	毫无疑问，seL4项目使用的验证器能力相当强。
	然而，seL4最初是一个单CPU项目。
	虽然seL4已经获得了多处理器能力，
	但它目前使用的是非常粗粒度的锁，
	类似于Linux内核的旧Big Kernel Lock (BKL)\@。
	希望有一天将seL4的验证器添加到一本关于并行编程的书中
	是有意义的，但现在还不是时候。
}\QuickQuizEnd

然而，直接从源代码或二进制文件进行验证都具有
消除人工翻译错误的优势，
这对于可靠的回归测试至关重要。

这并不是说具有专用语言的工具是无用的。
相反，它们对于设计时验证非常有帮助，
正如在\cref{chp:Formal Verification}中讨论的那样。
然而，这类工具对于自动化回归测试并不特别有帮助，
而这事实上正是本节的主题。

\subsection{环境}
\label{sec:future:Environment}

形式化验证工具正确建模其环境至关重要。
一个常见的遗漏是\IX{memory model}，
其中大量形式化验证工具（包括Promela/\co{spin}）
被限制为\IXalth{sequential consistency}{sequential}{memory consistency}。
\cref{sec:formal:Is QRCU Really Correct?}中叙述的QRCU经历
是一个重要的警示故事。

Promela和\co{spin}假设顺序一致性，
这与现代计算机系统不太匹配，
正如在\cref{chp:Advanced Synchronization: Memory Ordering}中所见。
相比之下，PPCMEM和\co{herd}的一大优势是
它们对各种CPU家族内存模型的详细建模，
包括x86、\ARM、Power，以及在\co{herd}的情况下，
还有Linux内核内存模型~\cite{Alglave:2018:FSC:3173162.3177156}，
该模型已被接受进入Linux内核v4.17版本。

\co{cbmc}和Nidhugg工具提供了一些选择内存模型的能力，
但没有提供PPCMEM和\co{herd}那样的多样性。
然而，随着时间推移，更大规模的工具可能会采用更多种类的内存模型。

从更长远来看，如果形式化验证工具能够包含
I/O~\cite{PaulEMcKenney2016LinuxKernelMMIO}会很有帮助，
但这可能还需要一段时间才能实现。

尽管如此，无法匹配环境的工具仍然可以是有用的。
例如，许多并发bug在神话般的顺序一致系统上仍然是bug，
这些bug可以被用顺序一致性过度近似系统内存模型的
工具所发现。
然而，这些工具将无法找到涉及缺少内存排序指令的bug，
正如前述\cref{sec:formal:Is QRCU Really Correct?}
的警示故事中所指出的那样。

\subsection{开销}
\label{sec:future:Overhead}

几乎所有硬核形式化验证工具本质上都是指数级的，
这可能看起来令人沮丧，直到你考虑到许多最有趣的
软件问题实际上是不可判定的。
然而，即使在指数级之间也存在程度差异。

PPCMEM在设计上是未优化的，
以便更好地保证所关注的内存模型被准确表示。
\co{herd}工具则更激进地优化，
如\cref{sec:formal:Axiomatic Approaches}中所述，
因此比PPCMEM快了数个数量级。
然而，PPCMEM和\co{herd}都针对的是非常小的litmus测试，
而不是较大的代码体。

相比之下，Promela/\co{spin}、\co{cbmc}和Nidhugg是为
（某种程度上）更大的代码体设计的。
Promela/\co{spin}被用于验证好奇号火星车的
文件系统~\cite{DBLP:journals/amai/GroceHHJX14}，
如前所述，\co{cbmc}和Nidhugg都被应用于Linux内核RCU\@。

如果启发式方法的进步继续保持过去三十年的速度，
我们可以期待形式化验证的开销大幅减少。
也就是说，组合爆炸仍然是组合爆炸，
这将预期会严格限制可以被验证的程序的大小，
无论启发式方法是否继续改进。

然而，组合爆炸的另一面是马其顿的腓力二世的
永恒建议：``分而治之。''
如果一个大程序可以被分割并验证各个部分，
结果可以是组合\emph{内爆}~\cite{PaulEMcKenney2011Verico}。
一个自然的分割点是在API边界上，例如锁原语的边界。
一次验证通过可以验证锁实现是正确的，
额外的验证通过可以验证锁API的正确使用。

\begin{listing}
\input{CodeSamples/formal/herd/C-SB+l-o-o-u+l-o-o-u-C@whole.fcv}
\caption{用\tco{cmpxchg_acquire()}模拟锁}
\label{lst:future:Emulating Locking with cmpxchg}
\end{listing}

\begin{table}
\rowcolors{1}{}{lightgray}
\renewcommand*{\arraystretch}{1.1}
\small
\centering
\begin{tabular}{S[table-format=1.0]S[table-format=1.3]S[table-format=2.3]}
	\toprule
	\multicolumn{1}{c}{\# 线程数} & \multicolumn{1}{c}{锁} &
			\multicolumn{1}{c}{\tco{cmpxchg_acquire}} \\
	\midrule
	2 & 0.004 &  0.022 \\
	3 & 0.041 &  0.743 \\
	4 & 0.374 & 59.565 \\
	5 & 4.905 &        \\
	\bottomrule
\end{tabular}
\caption{模拟锁：
			    性能（秒）}
\label{tab:future:Emulating Locking: Performance (s)}
\end{table}

这种方法的性能优势可以使用Linux内核
内存模型~\cite{Alglave:2018:FSC:3173162.3177156}来演示。
该模型提供了\co{spin_lock()}和\co{spin_unlock()}原语，
但这些原语也可以使用\co{cmpxchg_acquire()}和
\co{smp_store_release()}来模拟，
如\cref{lst:future:Emulating Locking with cmpxchg}所示
（\path{C-SB+l-o-o-u+l-o-o-*u.litmus}和\path{C-SB+l-o-o-u+l-o-o-u*-C.litmus}）。
\Cref{tab:future:Emulating Locking: Performance (s)}
比较了使用模型的\co{spin_lock()}和\co{spin_unlock()}
与列表中所示的模拟这些原语的性能和可扩展性。
差异是显著的：
在四个进程时，模型比模拟快两个数量级以上！

\QuickQuiz{
\begin{fcvref}[ln:future:formalregress:C-SB+l-o-o-u+l-o-o-u-C:whole]
	为什么要在\cref{lst:future:Emulating Locking with cmpxchg}的
	\clnref{filter_}上使用单独的\co{filter}命令，
	而不是直接将条件添加到\co{exists}子句中？
	而且使用\co{xchg_acquire()}代替\co{cmpxchg_acquire()}
	不是更简单吗？
\end{fcvref}
}\QuickQuizAnswer{
	\co{filter}子句使得\co{herd}工具在比\co{exists}子句
	更早的处理阶段丢弃执行，这提供了显著的加速。

\begin{table}
\rowcolors{7}{lightgray}{}
\renewcommand*{\arraystretch}{1.1}
\small
\centering
\begin{tabular}{S[table-format=1.0]S[table-format=1.3]S[table-format=2.3]
		S[table-format=3.3]S[table-format=2.3]S[table-format=3.3]}
	\toprule
	& & \multicolumn{2}{c}{\tco{cmpxchg_acquire()}}
		& \multicolumn{2}{c}{\tco{xchg_acquire()}} \\
	\cmidrule(l){3-4} \cmidrule(l){5-6}
	\multicolumn{1}{c}{\#} & \multicolumn{1}{c}{锁}
		& \multicolumn{1}{c}{\tco{filter}}
			& \multicolumn{1}{c}{\tco{exists}}
				& \multicolumn{1}{c}{\tco{filter}}
					& \multicolumn{1}{c}{\tco{exists}} \\
	\cmidrule{1-1} \cmidrule(l){2-2} \cmidrule(l){3-4} \cmidrule(l){5-6}
	2 & 0.004 &  0.022 &   0.039 &  0.027 &  0.058 \\
	3 & 0.041 &  0.743 &   1.653 &  0.968 &  3.203 \\
	4 & 0.374 & 59.565 & 151.962 & 74.818 & 500.96 \\
	5 & 4.905 &        &         &        &        \\
	\bottomrule
\end{tabular}
\caption{模拟锁：
			    性能比较（秒）}
\label{tab:future:Emulating Locking: Performance Comparison (s)}
\end{table}

	至于\co{xchg_acquire()}，这个原子操作无论锁获取是否成功
	都会执行写操作，这意味着使用\co{xchg_acquire()}的模型
	将比使用\co{cmpxchg_acquire()}的模型有更多操作，
	后者在获取失败的情况下不会执行写操作。
	更多的写操作意味着更多的组合爆炸，
	如\cref{tab:future:Emulating Locking: Performance Comparison (s)}所示
	（\path{C-SB+l-o-o-u+l-o-o-*u.litmus}、
	\path{C-SB+l-o-o-u+l-o-o-u*-C.litmus}、
	\path{C-SB+l-o-o-u+l-o-o-u*-CE.litmus}、
	\path{C-SB+l-o-o-u+l-o-o-u*-X.litmus}和
	\path{C-SB+l-o-o-u+l-o-o-u*-XE.litmus}）。
	该表清楚地表明\co{cmpxchg_acquire()}
	优于\co{xchg_acquire()}，
	而使用\co{filter}子句优于使用\co{exists}子句。
}\QuickQuizEnd

如果工具能够自动将大程序分割、验证各个部分，
然后验证部分的组合，那当然是非常有用的。
在此之前，大程序的验证将需要大量的人工干预。
这种干预最好通过脚本来中介，
以便更可靠地在每个版本发布时进行重复验证，
最好最终以适合持续集成的方式进行。
而Facebook的Infer工具已经朝着这个方向迈出了重要步伐，
通过组合性和抽象化~\cite{SamBlackshear2018RacerD,DinoDistefano2019FBstaticAnalysis}。

无论如何，我们可以期待形式化验证能力随时间继续增强，
任何这样的增强都将反过来提高形式化验证对回归测试的适用性。

\subsection{定位Bug}
\label{sec:future:Locate Bugs}

任何有一定规模的软件工件都包含bug。
因此，一个只报告bug存在与否的形式化验证工具并不特别有用。
需要的是一个至少提供\emph{一些}关于bug位置和性质信息的工具。

\co{cbmc}的输出包含映射回源代码的回溯信息，
类似于Promela/\co{spin}的回溯，Nidhugg也是如此。
当然，这些回溯可能相当长，分析它们可能相当乏味。
然而，这样做通常比用传统方式定位bug要快得多、愉快得多。

此外，形式化验证工具的最简单测试之一是bug注入。
毕竟，不仅我们中的任何人都可以写
\co{printf("VERIFIED\\n")}，而且事实是
形式化验证工具的开发人员和我们其他人一样容易犯错。
因此，仅仅宣称存在bug的形式化验证工具从根本上来说
不太值得信赖，因为更难在实际代码上对其进行验证。

撇开这些不谈，编写形式化验证工具的人可以利用现有工具。
例如，一个旨在仅确定严重但罕见的bug是否存在的工具
可以利用二分法。
如果被测程序的旧版本不包含该bug，
但新版本包含，那么可以使用二分法快速定位引入bug的提交，
这可能是找到并修复bug的足够信息。
当然，这种策略对于常见bug不太有效，
因为在这种情况下二分法会失败，
因为所有提交都至少有一个常见bug的实例。

因此，许多形式化验证工具提供的执行跟踪将继续是有价值的，
特别是对于复杂且难以理解的bug。
此外，最近的工作应用了\emph{不正确性逻辑}形式化方法，
类似于用于完整正确性证明的传统Hoare逻辑，
但其唯一目的是查找bug~\cite{PeterWOHearn2019incorrectnessLogic}。

\subsection{最小脚手架}
\label{sec:future:Minimal Scaffolding}

在过去，形式化验证研究人员要求提供完整的规范，
软件将根据该规范进行验证。
不幸的是，一个数学上严格的规范可能比实际代码还要大，
规范的每一行与代码的每一行一样可能包含bug。
一项证明代码忠实实现了规范的形式化验证工作
将是两者之间bug对bug兼容性的证明，
这可能不太有帮助。

更糟糕的是，许多软件工件的需求本质上是经验性的，
包括Linux内核的RCU~\cite{PaulEMcKenney2015RCUreqts1,PaulEMcKenney2015RCUreqts2,PaulEMcKenney2015RCUreqts3}。\footnote{
	或者，用形式化验证的术语来说，Linux内核的RCU具有
	\emph{不完整的规范}。}
对于这种常见类型的软件，完整的规范是一种礼貌的虚构。
完整的规范对于硬件同样是虚构的，
这一点在2017年末的Meltdown和Spectre侧信道攻击
中得到了充分证明~\cite{JannHorn2018MeltdownSpectre}。

这种情况可能会让人对实际软件和硬件工件的形式化验证
完全丧失希望，但事实证明可以做很多事情。
例如，设计和编码规则可以作为部分规范，
代码中包含的断言也是如此。
事实上，\co{cbmc}和Nidhugg等形式化验证工具都会检查
可以触发的断言，隐式地将这些断言作为规范的一部分。
然而，断言也是代码的一部分，这使得它们不太可能变得过时，
特别是如果代码也经过了压力测试。\footnote{
	你\emph{确实}对你的代码进行压力测试，对吧？}
\co{cbmc}工具还检查数组越界引用，
从而隐式地将其添加到规范中。
前述的不正确性逻辑也可以被认为使用了一个隐式的
无bug存在的规范~\cite{PeterWOHearn2019incorrectnessLogic}。

这种隐式规范方法很有道理，特别是如果你将形式化验证
不是视为完整的正确性证明，
而是视为一种具有与常见情况（即测试）不同优势和劣势的
替代验证形式。
从这个角度来看，软件总会有bug，因此任何有助于
找到这些bug的任何种类的工具都是非常好的。

\subsection{相关的Bug}
\label{sec:future:Relevant Bugs}

找到bug并修复它们当然是任何类型验证工作的全部意义。
显然，应避免误报。
但即使没有误报，bug和bug也是不同的。

例如，假设一个软件工件恰好有100个剩余bug，
每个bug平均每一百万年的运行时间才会出现一次。
进一步假设一个全知的形式化验证工具定位了所有100个bug，
开发人员勤勉地修复了它们。
这个软件工件的可靠性会发生什么变化？

答案是可靠性\emph{降低}了。

要理解这一点，请记住历史经验表明大约7\pct{}的修复
会引入新的bug~\cite{RexBlack2012SQA}。
因此，修复这100个bug（其组合平均故障间隔时间（MTBF）约为10,000年）
将引入另外七个bug。
历史统计表明每个新bug的MTBF远小于70,000年。
这反过来表明这七个新bug的组合MTBF
很可能远小于10,000年，
这又意味着善意修复原来100个bug的行为实际上降低了
整体软件的可靠性。

\QuickQuizSeries{%
\QuickQuizB{
	我们怎么知道已知bug的MTBF是对尚未发现的bug的
	MTBF的良好估计？
}\QuickQuizAnswerB{
	我们不知道，但这并不重要。

	要理解这一点，请注意7\pct{}的数字仅适用于
	随后被发现的注入bug：
	它必然忽略了任何从未被发现的注入bug。
	因此，已知bug的MTBF统计可能是随后被发现的注入bug的
	良好近似。

	本节的一个关键点是，我们应该更关心那些给用户带来不便的bug，
	而不是那些从未实际出现的其他bug。
	当然，这\emph{并非}说我们应该完全忽略尚未给用户带来
	不便的bug，只是我们应该适当地确定优先级，
	以便首先修复最重要和最紧急的bug。
}\QuickQuizEndB
%
\QuickQuizE{
	但形式化验证工具应该立即找到修复引入的所有bug，
	那么这为什么是个问题？
}\QuickQuizAnswerE{
	这是个问题，因为现实世界的形式化验证工具
	（与那些仅存在于形式化验证更激进倡导者的想象中的工具不同）
	不是全知的，因此只能定位某些类型的bug。
	仅举一个例子，形式化验证工具不太可能发现
	与遗漏断言对应的bug，或者等价地，
	与规范中未发现部分对应的bug。
}\QuickQuizEndE
}

更糟糕的是，想象另一个软件工件有一个bug平均每天失败一次，
还有99个每百万年失败一次的bug。
假设一个形式化验证工具定位了99个百万年的bug，
但未能找到那个一天一次的bug。
修复定位到的99个bug将花费时间和精力，降低可靠性，
而对那个可能正在造成尴尬甚至更糟糕后果的每天失败
毫无帮助。

因此，最好有一个能优先定位最麻烦bug的验证工具。
然而，如\cref{sec:future:Locate Bugs}中所述，
可以利用额外的工具。
一个强大的工具就是普通的测试。
了解了bug之后，应该可以为其构建特定的测试，
也可能使用\cref{sec:debugging:Hunting Heisenbugs}中描述的
一些技术来增加bug出现的概率。
这些技术应该允许计算bug原始故障率的粗略估计，
这反过来可以用于确定bug修复工作的优先级。

\QuickQuiz{
	但许多形式化验证工具一次只能找到一个bug，
	因此必须在工具能定位下一个bug之前修复每个bug。
	在这样的工具下如何确定bug修复工作的优先级？
}\QuickQuizAnswer{
	一种方法是提供一个可能不适合生产环境的简单修复，
	但允许工具定位下一个bug。
	另一种方法是限制配置或输入，
	使得到目前为止定位的bug不会发生。
	有许多类似的方法，但共同的主题是
	从工具的角度修复bug通常比构建和验证生产级别修复
	容易得多，关键是要优先考虑构建和验证
	生产级别修复所需的更大工作量。
}\QuickQuizEnd

最近有一些形式化验证工作优先考虑具有较少抢占的执行，
基于较少的抢占次数更可能发生的合理假设。

识别相关bug可能听起来要求太高了，
但如果我们要真正提高软件可靠性，这正是所需要的。

\subsection{形式化回归记分卡}
\label{sec:future:Formal Regression Scorecard}

\begin{table*}
% \rowcolors{6}{}{lightgray}
%\renewcommand*{\arraystretch}{1.1}
\small
\centering
\setlength{\tabcolsep}{2pt}
\begin{tabular}{lcccccccccc}
	\toprule
	& & Promela & & PPCMEM & & \tco{herd} & & \tco{cbmc} & & Nidhugg \\
	\midrule
	(1) 自动化 &
		& \cellcolor{red!50} &
			& \cellcolor{orange!50} &
				& \cellcolor{orange!50} &
					& \cellcolor{blue!50} &
						& \cellcolor{blue!50} \\
	\addlinespace[3pt]
	(2) 环境 &
		& \cellcolor{red!50} (MM) &
			& \cellcolor{green!50} &
				& \cellcolor{blue!50} &
					& \cellcolor{yellow!50} (MM) &
						& \cellcolor{orange!50} (MM) \\
	\addlinespace[3pt]
	(3) 开销 &
		& \cellcolor{yellow!50} &
			& \cellcolor{red!50} &
				& \cellcolor{yellow!50} &
					& \cellcolor{yellow!50} (SAT) &
						& \cellcolor{green!50} \\
	\addlinespace[3pt]
	(4) 定位Bug &
		& \cellcolor{yellow!50} &
			& \cellcolor{yellow!50} &
				& \cellcolor{yellow!50} &
					& \cellcolor{green!50} &
						& \cellcolor{green!50} \\
	\addlinespace[3pt]
	(5) 最小脚手架 &
		& \cellcolor{green!50} &
			& \cellcolor{yellow!50} &
				& \cellcolor{yellow!50} &
					& \cellcolor{blue!50} &
						& \cellcolor{blue!50} \\
	\addlinespace[3pt]
	(6) 相关Bug &
		& \cellcolor{yellow!50} ??? &
			& \cellcolor{yellow!50} ??? &
				& \cellcolor{yellow!50} ??? &
					& \cellcolor{yellow!50} ??? &
						& \cellcolor{yellow!50} ??? \\
	\bottomrule
\end{tabular}
\caption{形式化回归记分卡}
\label{tab:future:Formal Regression Scorecard}
\end{table*}

\Cref{tab:future:Formal Regression Scorecard}
显示了本章涵盖的形式化验证工具的粗略记分卡。
波长越短越好。

Promela需要手动翻译且仅支持
\IXalth{sequential consistency}{sequential}{memory consistency}，
因此其前两个单元格为红色。
它有合理的开销（就形式化验证而言）
且提供回溯信息，因此其接下来两个单元格为黄色。
尽管需要手动翻译，Promela以自然的方式处理断言，
因此其第五个单元格为绿色。

PPCMEM通常需要手动翻译，因为它支持的litmus测试规模很小，
因此其第一个单元格为橙色。
它处理多种\IX{memory model}，因此其第二个单元格为绿色。
它的开销相当高，因此其第三个单元格为红色。
它提供操作之间关系的图形显示，
虽不如回溯信息有帮助，但仍然相当有用，
因此其第四个单元格为黄色。
它需要构造\co{exists}子句且不能进行进程内断言，
因此其第五个单元格也是黄色。

\co{herd}工具具有与PPCMEM类似的大小限制，
因此\co{herd}的第一个单元格也是橙色。
它支持多种内存模型，因此其第二个单元格为蓝色。
它有合理的开销，因此其第三个单元格为黄色。
它的bug定位和断言能力与PPCMEM非常相似，
因此\co{herd}接下来两个单元格也是黄色。

\co{cbmc}工具直接输入C代码，因此其第一个单元格为蓝色。
它支持几种内存模型，因此其第二个单元格为黄色。
它有合理的开销，因此其第三个单元格也是黄色，
不过也许SAT求解器的性能将继续改善。
它提供回溯信息，因此其第四个单元格为绿色。
它直接从C代码中获取断言，因此其第五个单元格为蓝色。

Nidhugg也直接输入C代码，因此其第一个单元格也是蓝色。
它只支持几种内存模型，因此其第二个单元格为橙色。
它的开销相当低（就形式化验证而言），
因此其第三个单元格为绿色。
它提供回溯信息，因此其第四个单元格为绿色。
它直接从C代码中获取断言，因此其第五个单元格为蓝色。

那么第六行也就是最后一行呢？
现在判断这些工具在找到正确bug方面表现如何还为时过早，
因此它们都是带问号的黄色。

\QuickQuizSeries{%
\QuickQuizB{
	在\cref{tab:future:Formal Regression Scorecard}
	所示的记分卡中，测试会如何表现？
}\QuickQuizAnswerB{
	它将一路蓝色到底，
	可能的例外是第三行（开销），
	可能会因为测试难以找到不太可能出现的bug而被扣分。

	另一方面，不太可能出现的bug通常也是不相关的bug，
	因此你的结果可能会有所不同。

	很大程度上取决于你的安装基础的大小。
	如果你的代码只会在（比如说）10,000个系统上运行，
	Murphy实际上可以是一个非常好的人。
	所有可能出错的事情都会出错。
	最终。
	也许在地质时间尺度上。

	但如果你的代码运行在200亿个系统上，
	就像据说2017年底Linux内核那样，
	Murphy可以是一个真正的混蛋！
	所有可能出错的事情都会出错，
	而且可能出错得非常快！！！
}\QuickQuizEndB
%
\QuickQuizE{
	难道不是有比\cref{tab:future:Formal Regression Scorecard}
	中显示的更多的形式化验证系统吗？
}\QuickQuizAnswerE{
	确实如此！
	该表关注的是Paul使用过的那些，但其他一些也在证明其有用性。
	形式化验证在seL4项目中得到了大量
	使用~\cite{ThomasSewell2013L4binaryVerification}，
	其工具现在可以处理适度水平的并发。
	最近，Catalin Marinas使用Lamport的
	TLA工具~\cite{Lamport:2002:SST:579617}定位了Linux内核
	排队自旋锁（spinlock）实现中的一些前向进展bug。
	Will Deacon修复了这些bug~\cite{WillDeacon2018qspinlock}，
	而Catalin验证了Will的
	修复~\cite{CatalinMarinas2018qspinlockTLA}。

	较轻量级的形式化验证工具已在生产中得到大量
	使用~\cite{JamesRLarus2004RightingSoftware,AlBessey2010BillionLoCLater,ByronCook2018FormalAmazon,CaitlinSadowski2018staticAnalysisGoogle,DinoDistefano2019FBstaticAnalysis}。
}\QuickQuizEndE
}

再次请注意，该表对这些工具在回归测试中的使用进行评分。
仅仅因为它们中的许多不适合回归测试，
并不意味着它们是无用的，事实上，
它们中的许多已经多次证明了自己的价值。\footnote{
	仅举一个例子，Promela被用于验证好奇号火星车的文件系统。
	\emph{你的}形式化验证工具是否被用于目前在火星上运行的软件？？？}
只是不适合回归测试。

然而，这可能会改变。
毕竟，形式化验证工具在2010年代取得了令人印象深刻的进步。
如果这种进步继续下去，形式化验证可能会成为
并行程序员验证工具箱中不可或缺的工具。

% future/fp.tex
% mainfile: ../perfbook.tex
% SPDX-License-Identifier: CC-BY-SA-3.0

\section{用于并行化的函数式编程}
\label{sec:future:Functional Programming for Parallelism}
%
\epigraph{函数式编程在并行应用中的奇怪失败。}
	 {Malte Skarupke}

当我在1980年代初上第一门functional programming课时，
教授断言无副作用的functional programming风格非常适合
轻松实现并行化和分析。
三十年后，这一断言依然存在，但主流生产环境中
并行函数式语言的使用微乎其微——这种状况可能与教授的另一个断言
（即程序既不应该维护状态也不应该进行I/O）并非完全无关。
Erlang等函数式语言有小众应用，
多线程支持也已添加到其他几种函数式语言中，
但主流生产使用仍然是C、C++、Java和Fortran等过程式语言的领地
（通常与OpenMP、MPI或coarray配合使用）。

这种情况自然引出一个问题：``如果分析是目标，
为什么不在分析之前将过程式语言转换为函数式语言？''
当然，对这种方法有许多反对意见，其中我列出三条：

\begin{enumerate}
\item	过程式语言经常大量使用全局变量，
	这些变量可以被不同的函数独立更新，
	或者更糟糕的是，被多个线程独立更新。
	注意，Haskell的\emph{monad}就是为了处理
	单线程全局状态而发明的，而多线程访问
	全局状态对函数式模型造成了更多的伤害。
\item	多线程过程式语言经常使用锁、原子操作和事务等
	同步原语，这些对函数式模型造成了更多的伤害。
\item	过程式语言可以对函数参数进行\emph{别名化}，
	例如，通过将指向同一结构的指针作为两个不同的参数
	传递给同一函数的同一次调用。
	这可能导致函数在不知情的情况下通过两个不同的
	（可能重叠的）代码序列更新该结构，
	这极大地使分析变得复杂。
\end{enumerate}

当然，鉴于全局状态、同步原语和别名的重要性，
聪明的functional programming专家已经提出了许多尝试
来调和函数式编程（functional programming）模型与它们之间的矛盾，
monad只是其中一个例子。

另一种方法是将并行过程式程序编译为函数式程序，
然后使用functional programming工具来分析结果。
但完全可以比这做得更好，因为任何真实的计算
都是一个具有有限输入并运行有限时间的大型有限状态机。
这意味着任何真实程序都可以转换为一个表达式，
尽管可能是一个大得不切实际的表达式~\cite{VijayDSilva2012-sas}。

然而，并行算法的许多低级内核可以转换为
足够小的表达式，可以轻松装入现代计算机的内存中。
如果这样的表达式与一个断言相结合，检查该断言是否会触发
就变成了一个可满足性问题。
尽管可满足性问题是NP完全的，但通常可以在
比生成完整状态空间所需的时间短得多的时间内解决。
此外，求解时间似乎只弱依赖于底层内存模型，
因此运行在弱排序系统上的算法也可以被检查~\cite{JadeAlglave2013-cav}。

一般方法是将程序转换为静态单赋值（SSA）形式，
使得对变量的每次赋值都创建该变量的一个单独版本。
这适用于所有活跃线程的赋值，因此生成的表达式
包含了所讨论代码的所有可能执行。
添加一个断言就相当于询问输入和初始值的任何组合
是否可能导致断言触发，如上所述，
这正是可满足性问题。

一个可能的反对意见是它不能优雅地处理任意循环构造。
然而，在许多情况下，可以通过将循环展开有限次数来处理。
此外，也许某些循环也可以通过归纳方法来折叠。

另一个可能的反对意见是自旋锁涉及任意长的循环，
任何有限的展开都无法捕获自旋锁的完整行为。
事实证明这个反对意见很容易克服。
不是建模完整的自旋锁，而是建模一个trylock，它尝试
获取锁，如果不能立即获取就中止。
然后必须精心构造断言，以避免在自旋锁因锁不可立即获取
而中止的情况下触发。
因为逻辑表达式与时间无关，所有可能的并发行为
都将通过这种方法被捕获。

最后一个反对意见是，这种技术不太可能能够处理
像Linux内核这样数百万行代码构成的完整规模的软件工件。
这很可能是事实，但事实仍然是，对Linux内核中
每个小得多的并行原语进行穷举验证将是相当有价值的。
事实上，推动这一方法的研究人员已经将其应用于
非平凡的实际代码，包括Linux内核中的Tree RCU
实现~\cite{LihaoLiang2016VerifyTreeRCU,MichalisKokologiannakis2017NidhuggRCU}。

这种技术的适用范围有多广还有待观察，但它是
形式化验证领域中更有趣的创新之一。
虽然functional programming的倡导者终于可能在其
函数式编程必将统治一切的断言上是正确的，但显然
这种长期被吹捧的方法论正在其形式化验证的主场
开始面临可信的竞争。
因此，仍然有理由怀疑functional programming统治的不可避免性。

% @@@ Maybe add quantum computing back in, but heavily summarized.

\section{总结}
\label{sec:future:Summary}
%
\epigraph{生活只能在回顾中理解；
	   但必须向前过。}
	  {S{\o}ren Kierkegaard}

本章快速浏览了若干可能的未来，
包括多核、transactional memory、作为回归测试的形式化验证，
以及并发functional programming。
这些未来中的任何一个都可能成为现实，但更有可能的是，
与过去一样，未来将比我们所能想象的更加奇异。

\QuickQuizAnswersChp{qqzfuture}
